% =====================================================================
%  Testing, Experimentation and Validation  (~3 pages)
%  A SECTION of the FYP report. Include from the parent chapter with:
%      % =====================================================================
%  Testing, Experimentation and Validation  (~3 pages)
%  A SECTION of the FYP report. Include from the parent chapter with:
%      % =====================================================================
%  Testing, Experimentation and Validation  (~3 pages)
%  A SECTION of the FYP report. Include from the parent chapter with:
%      % =====================================================================
%  Testing, Experimentation and Validation  (~3 pages)
%  A SECTION of the FYP report. Include from the parent chapter with:
%      \input{testing_validation}
%
%  Intended to follow \input{hitl_edge_processing}; cross-references
%  \ref{sec:hitl}. A longer version with the full 26-case test matrix is
%  kept in testing_validation_long.tex if an appendix is wanted.
%
%  REQUIRED PACKAGES
%      \usepackage{booktabs}
%      \usepackage{array}
%      \usepackage{graphicx}
%      \usepackage{subcaption}
%      \graphicspath{{figures/}}
%
%  IMAGES EXPECTED in docs/report/figures/  (see checklist at end)
%      detection_a.png  detection_b.png  detection_c.png
% =====================================================================

\section{Testing, Experimentation and Validation}
\label{sec:test}

Testing of the hardware-in-the-loop testbed described in
Section~\ref{sec:hitl} was organised around a difficulty specific to this
kind of system: a fault at a low layer does not raise an error at the layer
above it, but produces plausible and wrong behaviour there. A subscriber
starved of buffer space receives nothing and reports nothing; two nodes with
mismatched topic names run indefinitely without complaint. Testing was
therefore performed bottom-up --- physical link, path separation, middleware
transport, detector, end-to-end --- with each layer required to pass before
the next was attempted.

Three results are reported here. The first establishes that the two-link
architecture does what it claims. The second establishes that the detector
finds real people rather than plausible-looking texture. The third concerns
the measurement chain itself, and is included because it determines whether
the other two can be believed.

% ---------------------------------------------------------------------
\subsection{Test environment}
\label{sec:test-env}

Both machines run Ubuntu~22.04 with ROS\,2 Humble and Fast~DDS~2.6, so the
middleware is identical at each end of the link and cannot itself be a
source of difference. The simulation host runs Gazebo Classic in lockstep
with ArduPilot SITL, and ns-3 with TapBridge for the wireless channel. The
edge node is a Raspberry~Pi~4B running YOLOv8n restricted to the COCO
\emph{person} class. The camera produces $1280 \times 720$ frames of
$2.76$\,MB at $5$\,Hz. The test scene contains exactly five human models at
coordinates fixed in the world definition, which makes recall a counted
quantity rather than an estimated one.

The Pi's single Ethernet port is divided by 802.1Q tagging into
\texttt{eth0.10} ($10.0.0.2$, camera, unimpaired) and \texttt{eth0.42}
($10.42.0.12$, detections, through ns-3), created by
\texttt{scripts/netns/pi\_hitl\_link.sh}. The matching host interfaces
\texttt{eth-cam} and \texttt{eth-rf} are created by step~1c of
\texttt{scripts/netns/launch\_single\_uav\_netns.sh}, the latter bridged
into \texttt{tap-uav2}.

% ---------------------------------------------------------------------
\subsection{Result 1: the two links behave differently}
\label{sec:test-links}

The central claim of the architecture is that detection results cross a
simulated wireless channel while camera imagery does not. Two paths that
both work are not evidence of this; they must also be shown to differ in the
intended way, and imagery must be shown to be absent from the radio path.

Table~\ref{tab:test-links} reports sixty ICMP samples on the sensor path and
thirty on the radio path. Large samples are necessary: the first packet on the
radio path includes address resolution across the simulated channel, and a
two-sample measurement of the same configuration returned a mean of
$167$\,ms against $74.5$\,ms for thirty. The short sample is biased, not
merely noisy.

\begin{table}[htbp]
\centering
\caption{Measured characteristics of the two paths. Sensor path over sixty
samples, radio path over thirty; zero packet loss on both.}
\label{tab:test-links}
\begin{tabular}{@{}lrrrr@{}}
\toprule
Path & Min (ms) & Mean (ms) & Max (ms) & Jitter (ms) \\
\midrule
Sensor, VLAN 10 (wired)   & 1.0  & 1.4  & 1.8   & 0.12 \\
Radio, VLAN 42 (via ns-3) & 10.1 & 74.5 & 158.0 & 38.4 \\
\midrule
Ratio                     & --   & $53\times$ & -- & $320\times$ \\
\bottomrule
\end{tabular}
\end{table}

The radio path is approximately \(53\times\) slower on average, and its
variation is larger by more than two orders of magnitude.

Its internal spread is itself informative: a minimum of \(10.1\)\,ms against
a mean of \(74.5\)\,ms is a range of some \(7\times\) on one link, where a
channel applying a fixed delay would produce a narrow distribution. The
variation indicates that ns-3 is modelling contention, in which packets queue
and back off, rather than adding a constant offset.

The radio measurement was repeated. An independent set of thirty samples taken
several hours earlier, with the host reconfigured in between, gave a mean of
\(77.4\)\,ms against the \(74.5\)\,ms reported here, agreeing within
\(4\%\).

The sensor figures were taken \emph{while} the camera was streaming at
\(110\)\,Mbps. An idle-cable measurement over the same number of samples
gives the same mean, so the two logical links sharing one physical cable do
not delay one another. Queue management was verified rather than assumed:
both machines run \texttt{fq\_codel} with a \(5\)\,ms target, and byte
queue limits on the edge node had reduced its driver ring to approximately one
packet.

Absence of imagery on the radio path was established negatively, by
observation rather than by inspecting a rule set, since separation here is
enforced by addressing and bridge membership rather than by filtering. No
$10.0.0.0/24$ traffic was observed on \texttt{tap-uav2}, the point at which
traffic enters ns-3, and the byte counters on the two interfaces differ by
approximately four orders of magnitude --- consistent with $2.76$\,MB per
frame on one path against $118$\,B per detection message on the other, a
reduction of approximately \(23{,}000\times\).

A structural argument supports the measurements. The ground station receiver
executes inside the \texttt{gcsns} network namespace, whose only interface
is one end of a veth pair attached to \texttt{tap-gcs}. That namespace has
no interface on $10.0.0.0/24$ and no route to it. A message arriving at the
receiver cannot have taken any other path, because no other path exists from
inside the namespace. Separation is therefore a property of the topology
rather than of a setting that could be silently overridden.

% ---------------------------------------------------------------------
\subsection{Result 2: detections correspond to real subjects}
\label{sec:test-detection}

Visual confirmation was made a precondition of accepting any quantitative
result. A detector reporting a plausible count while placing its boxes on
building texture would satisfy every numerical check in this section, and an
earlier configuration did exactly that: sixteen confident detections against
five subjects, none correct. Only inspection of annotated frames
distinguishes the two cases. Those false positives were eliminated by
raising the confidence threshold and narrowing the camera field of view.

\begin{figure}[htbp]
\centering
\begin{subfigure}[b]{0.32\textwidth}
  \includegraphics[width=\textwidth]{detection_a.png}
  \caption{Two subjects, confidence $0.59$ and $0.42$.}
  \label{fig:test-det-a}
\end{subfigure}
\hfill
\begin{subfigure}[b]{0.32\textwidth}
  \includegraphics[width=\textwidth]{detection_b.png}
  \caption{One subject, confidence $0.48$.}
  \label{fig:test-det-b}
\end{subfigure}
\hfill
\begin{subfigure}[b]{0.32\textwidth}
  \includegraphics[width=\textwidth]{detection_c.png}
  \caption{Three subjects, confidence $0.59$, $0.47$ and $0.46$.}
  \label{fig:test-det-c}
\end{subfigure}
\caption{Detector output during flight, with the edge node's log above each
view. Bounding boxes coincide with the human models rather than with ground
texture. The log records every frame, including those with no detection, so
that a stalled node is distinguishable from an empty scene, and reports
inference duration per frame.}
\label{fig:test-detection}
\end{figure}

Figure~\ref{fig:test-detection} shows three moments from a patrol flight.
Between one and four of the five subjects were detected per frame, at
confidences of \(0.42\)--\(0.59\). The shortfall against the known count of
five is attributable to scene geometry rather than to detector deficiency:
the models stand approximately one metre apart, so from the oblique camera
view the nearer figures partially occlude those behind them, and a partially
occluded subject at approximately \(53\) pixels in height falls below the
confidence threshold. Because the number of subjects is known exactly, any
count exceeding five would be a false positive without the need to
adjudicate individual boxes; none was observed in this configuration.

Inference duration was measured around the model invocation alone, excluding
message reception, format conversion and publication, so that the figure
remains comparable between back-ends and against standalone execution. Over
\(480\) consecutive frames of one patrol flight the mean was
\(1{,}093\)\,ms, with a minimum of \(1{,}025\)\,ms. The distribution is tight
apart from a single outlier at \(6{,}949\)\,ms, one frame in \(480\),
attributable to contention for the processor from another task rather than to
the model itself.

The Pi therefore completes rather less than one detection per second against a
camera delivering five frames per second, and discards approximately \(80\%\)
of what it receives --- which is why the camera subscription uses
\texttt{BEST\_EFFORT} with queue depth one, keeping the newest frame rather
than accumulating a backlog of stale ones.

% ---------------------------------------------------------------------
\subsection{Result 3: validating the measurement chain}
\label{sec:test-chain}

The two results above depend on the camera actually reaching the edge node
at the rate the configuration specifies. That assumption failed, and the
manner of its failure is reported here because it determines what the other
measurements are worth.

The symptom was that the edge node received imagery at \(0.19\)\,Hz where
the camera was configured for \(5\)\,Hz, a shortfall of approximately
\(26\times\). The visible consequences appeared three layers above the
cause: the detector reported no humans in \(93\) of \(94\) frames, and the
patrol mission timed out at every waypoint. Both invite explanation in terms
of detector settings or flight control, and neither has anything to do with
either --- the aircraft flew correctly and the detector functioned
correctly, but the scene was being observed through roughly one frame every
five seconds.

\begin{table}[htbp]
\centering
\caption{Eliminating candidate causes of the frame-rate shortfall. Each row
is a measurement that excludes one layer.}
\label{tab:test-elimination}
\begin{tabular}{@{}p{38mm}p{38mm}p{40mm}@{}}
\toprule
Hypothesis & Measurement & Outcome \\
\midrule
Simulation running slowly
  & Real-time factor from \texttt{/clock}
  & \(0.998\); simulated time tracks wall clock \\[3pt]
Gazebo rendering slowly
  & Host-side subscriber on the camera topic
  & \(5.00\)\,Hz, gaps \(0.19\)--\(0.21\)\,s; source correct \\[3pt]
Link negotiated at $100$\,Mbps
  & \texttt{ethtool} on the edge node
  & $1000$\,Mb/s full duplex \\[3pt]
Fragments lost in flight
  & Interface counters, both machines
  & Zero errors, zero drops \\[3pt]
Receive buffers too small
  & \texttt{net.core.rmem\_max}
  & $512$\,MB, already raised \\[3pt]
\textbf{Send buffers too small}
  & \texttt{net.core.wmem\_max}
  & \(\mathbf{208}\)\,\textbf{kB}, left at the default \\
\bottomrule
\end{tabular}
\end{table}

The DDS profile requests a \(16\)\,MB send buffer. The kernel caps any such
request at \texttt{net.core.wmem\_max}, and performs the reduction silently.
A single frame is therefore approximately \(13\times\) larger than the
buffer available to transmit it: the publisher filled that buffer
immediately and discarded the remaining fragments at the socket layer,
before they reached the network interface. Comparing production with
transmission confirmed this --- the host generated \(414\)\,MB of imagery
over thirty seconds and placed \(10\)\,MB of it on the wire, \(2\)\,Mbps
against the \(110\)\,Mbps the configuration implies. Raising the parameter
on both machines and restarting the stack, which is required because
Fast~DDS applies socket options only at participant creation, restored
delivery to \(5.000\)\,Hz with an inter-arrival standard deviation of
\(4.5\)\,ms: an exact match to the source.

Two properties of this fault are worth recording. First, no counter
registers it. Loss at the socket layer is invisible to interface
statistics, and \texttt{tx\_dropped} on the sending interface and
\texttt{rx\_errors} on the receiving one both remained at zero throughout,
because the discarded fragments were never transmitted. A diagnosis founded
on interface counters would have concluded that the network was healthy,
which it was. Second, the two directions are governed by separate kernel
parameters. The receive path had been corrected during earlier work and the
transmit path had not, because the failure that prompted the original
correction pointed only at the receiving side. The pre-flight checks in
\texttt{scripts/netns/run\_hitl.sh} now verify both parameters on both
machines before the stack starts.

% ---------------------------------------------------------------------
\subsection{Limitations}
\label{sec:test-limitations}

Clock synchronisation between the two machines has not been established.
End-to-end latency would be computed from timestamps taken on different
hosts, and without synchronisation such values are not meaningful while
remaining plausible in magnitude. End-to-end latency is therefore not
reported, rather than reported with a caveat.

The edge configuration is characterised; the corresponding ground
configuration, in which compressed imagery crosses the simulated channel and
inference is performed at the ground station, remains to be measured. All
measurements were taken with the aircraft close to the ground station,
giving zero packet loss on both paths, so the reliability policy applied to
detection messages is untested under the loss conditions it exists to
handle. The testbed supports a single physical edge node; multi-node
operation has not been evaluated.

% =====================================================================
%  IMAGE CHECKLIST — save into docs/report/figures/
%
%  detection_a.png / detection_b.png / detection_c.png
%      The three screenshots already captured: the detector log above and
%      the rqt_image_view annotated frame below, showing 2, 1 and 3
%      detections respectively.
%
%      CROP EACH ONE before including. At three-to-a-row the full desktop
%      screenshot will be illegible. Keep only:
%        - the last ~8 log lines (enough to show the count and the
%          inference time), and
%        - the image pane itself, without the rqt toolbar and topic
%          selector.
%      Crop all three to the same aspect ratio so the row lines up.
%
%      If the log text is still too small at 0.32\textwidth, drop the log
%      from the figure entirely and show only the three annotated views —
%      the inference figures are already stated in the prose.
%
%  OPTIONAL, if a fourth image is wanted:
%      testbed_hardware.jpg — a photograph of the laptop, the USB gigabit
%      adapter, the single Ethernet cable and the Pi 4B. This is the only
%      image that shows the work is hardware-in-the-loop rather than pure
%      simulation. It would go in Section \ref{sec:test-env}.
% =====================================================================

%
%  Intended to follow % =====================================================================
%  Hardware-in-the-Loop Edge Processing
%  A SECTION of the FYP report. Include from the parent chapter with:
%      \input{hitl_edge_processing}
%
%  HEADING LEVELS USED
%      \section     - the HITL section itself
%      \subsection  - its parts
%      \paragraph   - short findings within Results
%
%  REQUIRED PACKAGES (add to the main preamble if not already present)
%      \usepackage{booktabs}
%      \usepackage{graphicx}
%      \usepackage{tikz}
%      \usetikzlibrary{arrows.meta, positioning, fit, backgrounds}
%
%  All labels are prefixed hitl / tab:hitl- / fig:hitl- so they cannot
%  collide with labels elsewhere in the shared report.
% =====================================================================

\section{Hardware-in-the-Loop Edge Processing}
\label{sec:hitl}

The vision workload was moved off the simulation host and onto physical
hardware, so that the cost of onboard detection could be measured on a processor
representative of one that would actually fly. A Raspberry~Pi~4B performs
detection on imagery produced by the simulation, and returns its results across a
wireless channel simulated in ns-3.

% ---------------------------------------------------------------------
\subsection{Extending the testbed with physical hardware}
\label{sec:hitl-testbed}

The existing stack ran entirely on one machine: Gazebo for the environment and
sensors, ArduPilot SITL for the autopilot, ns-3 for the wireless channel, and
Linux network namespaces isolating the ground station from each aircraft. The
ns-3 scenario defined four nodes --- a ground control station and three aircraft
--- of which the second and third were unused placeholders.

A Raspberry~Pi~4B was introduced as the physical edge node and assigned to the
second aircraft slot. Its TAP device, previously created only to satisfy the
scenario's argument list, was brought into service as the attachment point. The
autopilot, physics and channel simulation remained on the host; only the vision
computation moved onto hardware.

The Pi was configured with Ubuntu~22.04 and ROS\,2 Humble to match the host
exactly, giving both machines the same DDS implementation. The host has no
Ethernet port, so a USB~3.0 gigabit adapter was fitted. An initial low-cost
adapter was replaced after \texttt{iperf3} measured \(4.5\)\,Mbps with \(24\%\)
loss, despite \texttt{ethtool} reporting \(100\)\,Mbps full duplex; the
replacement sustained \(834\)\,Mbps with no retransmissions.

% ---------------------------------------------------------------------
\subsection{Constructing the two-link network}
\label{sec:hitl-network}

Two logically separate paths were required between host and Pi. One carries
camera imagery and must remain unimpaired, since on a real aircraft it is a short
cable inside a single airframe rather than a radio. The other carries detection
results and must pass through the simulated channel, as that is the link the
experiment exists to characterise. The same reasoning applies to the
flight-dynamics traffic between SITL and Gazebo, which also bypasses ns-3.

The Pi provides a single Ethernet interface, so both paths were carried over one
cable using IEEE~802.1Q VLAN tagging. On the host, two virtual interfaces were
created on the physical adapter: \texttt{eth-cam}, tagged VLAN~10 and addressed
\(10.0.0.1\); and \texttt{eth-rf}, tagged VLAN~42, left without an address and
made a member of a new bridge, \texttt{br-uav2}. That bridge also received
\texttt{tap-uav2}, the ns-3 interface for the second aircraft node. On the Pi,
matching interfaces \texttt{eth0.10} (\(10.0.0.2\)) and \texttt{eth0.42}
(\(10.42.0.12\)) were created.

Attaching the physical interface to the bridge directly was rejected during
design. A Linux bridge floods frames to all member ports, so camera traffic would
have been delivered into the ns-3 TAP device as well, injecting approximately
\(111\)\,Mbps of imagery into the simulated channel.

The resulting arrangement enforces separation through routing rather than
filtering. Gazebo is addressed only on \(10.0.0.1\) and can reach the Pi only at
\(10.0.0.2\); the ground station namespace is addressed only on \(10.42.0.10\)
and can reach it only at \(10.42.0.12\). Neither can use the other's path.

\begin{figure}[htbp]
\centering
\begin{tikzpicture}[
  font=\small,
  box/.style={draw, rounded corners=2pt, align=center, inner sep=5pt},
  proc/.style={box, fill=black!4},
  net/.style={box, fill=black!10},
  arr/.style={-{Latex[length=2mm]}, thick},
]
\node[proc] (gz)   {Gazebo\\camera + physics};
\node[proc, below=14mm of gz] (gcs) {Ground station\\(namespace)\\\texttt{10.42.0.10}};
\node[net,  right=16mm of gcs] (ns3) {ns-3\\channel};
\node[net,  above=8mm of ns3]  (br)  {\texttt{br-uav2}\\+ TAP};
\node[proc, right=46mm of gz] (pi) {Raspberry Pi 4B\\YOLOv8n detector};

\draw[arr] (gz) -- node[above, align=center, font=\scriptsize]
  {VLAN 10 \quad camera, 2.76\,MB/frame\\unimpaired} (pi);
\draw[arr] (pi.south) |- node[pos=0.75, above, font=\scriptsize]
  {VLAN 42 \quad detections, 118\,B} (br.east);
\draw[arr] (br) -- (ns3);
\draw[arr] (ns3) -- (gcs);

\begin{scope}[on background layer]
  \node[draw, dashed, rounded corners, fit=(gz)(gcs)(ns3)(br),
        inner sep=7pt, label={[font=\scriptsize\itshape]above left:Host PC}] {};
\end{scope}
\end{tikzpicture}
\caption{The two-link arrangement. Camera imagery reaches the edge node over an
unimpaired wired path; detection results traverse the simulated channel. Both
share one physical cable, separated by VLAN tagging.}
\label{fig:hitl-architecture}
\end{figure}

Addressing was made persistent using NetworkManager on the host and netplan on
the Pi, the two distributions using different network managers. Addresses
assigned imperatively were found to be removed within minutes, as each manager
reconciles kernel state against its own configuration and deletes entries it did
not create.

% ---------------------------------------------------------------------
\subsection{Configuring the middleware}
\label{sec:hitl-dds}

A Fast DDS profile was written and applied on both machines. It pins participants
to the intended interfaces through an interface whitelist, since both machines
are attached to WiFi as well as the wired link and DDS otherwise advertises on
every available interface. The profile also disables the shared-memory transport,
which is not isolated by network namespaces and would allow participants to
exchange data without traversing the simulated channel.

Kernel socket buffers were raised on both machines from the \(208\)\,kB default
to \(512\)\,MB. A single \(1280 \times 720\) frame is \(2.76\)\,MB and fragments
into approximately \(1{,}900\) UDP datagrams, of which the default buffer holds
about \(7\%\); reassembly never completed and subscribers received nothing, with
no error reported.

The camera subscription on the Pi was changed from the ROS\,2 default of
\texttt{RELIABLE} with queue depth ten to \texttt{BEST\_EFFORT} with depth one.
Under the default policy the publisher retains each sample until the subscriber
acknowledges it, and because the Pi acknowledges at approximately \(1\)\,Hz, the
camera output of the entire simulation was throttled to \(0.27\)\,Hz.
\texttt{RELIABLE} was retained for detection messages, which must not be lost.

% ---------------------------------------------------------------------
\subsection{Deploying the detector}
\label{sec:hitl-detector}

The detection node runs YOLOv8n restricted to the COCO \emph{person} class,
subscribing to the camera topic and publishing results as JSON. It was deployed
to the Pi and instrumented to record inference duration, measured around the
model invocation alone so that figures remain comparable between back-ends and
against standalone execution.

Two dependency conflicts were resolved during deployment. Installing
\texttt{ultralytics} brings in NumPy~2.x, against which \texttt{cv\_bridge} --- a
compiled extension built for the 1.x C ABI --- fails. Installing an older OpenCV
to compensate reintroduces NumPy~2 as a dependency, and OpenCV~5.0 separately
breaks \texttt{cv\_bridge} by renaming the type constants it uses to build its
conversion table. Versions were pinned to \texttt{numpy==1.26.4} and
\texttt{opencv-python==4.10.0.84}, the only combination satisfying both.

The model was additionally exported to NCNN for ARM-optimised inference. Because
NCNN fixes the input tensor shape at export time, the export geometry was
specified explicitly as \(544 \times 960\) to match the letterboxed frame; a
default export produces a square input in which approximately \(43\%\) of the
computation is spent on padding, at \(19.9\)\,GFLOPS against \(11.3\)\,GFLOPS for
the correct shape. Half-precision export was avoided, the Cortex-A72 lacking fp16
SIMD support.

% ---------------------------------------------------------------------
\subsection{Automation}
\label{sec:hitl-automation}

Bring-up was consolidated into a single script that starts the namespaces, ns-3,
Gazebo, SITL, the micro-ROS agent and the ground station receiver, creates the
VLANs and bridge, and launches the detector on the Pi over SSH. Pre-flight checks
were added for each failure encountered during development: presence of the host
address before Gazebo starts, since Fast DDS enumerates interfaces only at
participant creation and cannot use one assigned later; reachability of the Pi;
adequate socket buffer limits on both machines; and confirmation after startup
that Gazebo bound the intended interface.

% ---------------------------------------------------------------------
\subsection{Results}
\label{sec:hitl-results}

Detections were verified visually, by inspecting annotated frames, before any
quantitative result was accepted. An earlier configuration had reported up to
sixteen detections against five subjects; these were false positives on building
texture, eliminated by raising the confidence threshold and narrowing the camera
field of view.

\paragraph{Channel characterisation.}
Table~\ref{tab:hitl-links} compares the two paths, each measured over thirty ICMP
samples. Short samples are unreliable: the first packet includes address
resolution across the simulated channel, and a two-sample measurement of the same
configuration returned a mean of \(167\)\,ms against \(59\)\,ms for thirty.

\begin{table}[htbp]
\centering
\caption{Measured characteristics of the two links, thirty samples each.}
\label{tab:hitl-links}
\begin{tabular}{lrrrr}
\toprule
Path & Min (ms) & Mean (ms) & Max (ms) & Jitter (ms) \\
\midrule
Sensor, VLAN 10 (wired)   & 0.4  & 1.5  & 2.7   & 0.4  \\
Radio, VLAN 42 (via ns-3) & 11.7 & 77.4 & 242.6 & 54.8 \\
\bottomrule
\end{tabular}
\end{table}

Packet loss was zero on both paths, the aircraft remaining close to the ground
station throughout. The radio path is approximately \(53\times\) slower than the
wired path on average and varies approximately \(140\times\) more.

The distribution within the radio path is itself informative. Its minimum of
\(11.7\)\,ms against a mean of \(77.4\)\,ms is a spread of some \(6.6\times\) on
one link, where a channel applying a fixed delay would produce a narrow one.
The variation indicates that ns-3 is modelling contention, in which packets
queue and back off before transmission, rather than simply adding a constant
offset to each packet.

\paragraph{Sensor transport.}
With the camera at \(20\)\,Hz the wired path carried \(481\)\,Mbps at
\(42{,}544\) packets per second with zero dropped packets, occupying
approximately \(58\%\) of measured link capacity.

\paragraph{Inference latency.}
Table~\ref{tab:hitl-inference} reports latency for each combination of back-end
and input resolution. Changing back-end at constant resolution gives
approximately \(2.5\times\); halving the pixel count at constant back-end gives
approximately \(2.3\times\). Latency therefore scales close to linearly with
pixel count, a pixel ratio of \(0.47\) producing a time ratio of \(0.43\).

\begin{table}[htbp]
\centering
\caption{Inference latency on the Raspberry Pi 4B.}
\label{tab:hitl-inference}
\begin{tabular}{llrr}
\toprule
Back-end & Input & Latency (ms) & Rate (fps) \\
\midrule
PyTorch & $960 \times 544$ & 2540 & 0.39 \\
PyTorch & $640 \times 384$ & 1100 & 0.91 \\
NCNN    & $960 \times 544$ & 1000 & 1.00 \\
NCNN    & $640 \times 384$ & \multicolumn{2}{l}{not measured (est.\ 470)} \\
\bottomrule
\end{tabular}
\end{table}

\paragraph{Detection performance.}
Between two and four of the five subjects were detected per frame, a recall of
\(40\)--\(80\%\). The shortfall is attributable to scene geometry rather than
detector deficiency: the models are separated by approximately one metre, so from
the oblique camera view the nearer figures partially occlude those behind them,
and partially occluded subjects at approximately \(53\) pixels in height fall
below the confidence threshold. Because the number and position of subjects are
known exactly from the world definition, recall is a directly measurable quantity
rather than an estimate.

\paragraph{Transmitted volume.}
Each detection message is approximately \(118\)\,B, against \(2.76\)\,MB for the
corresponding frame, a reduction of approximately \(23{,}000\times\) in data
crossing the radio.

% ---------------------------------------------------------------------
\subsection{Contention between sensor transport and inference}
\label{sec:hitl-contention}

The two workloads on the edge node proved not to be independent. The same model,
on the same input, required \(623\)\,ms executed on a static image but
\(1{,}343\)\,ms within the live pipeline: a difference of \(720\)\,ms, or
approximately half the available inference throughput.

The edge node received frames at \(19.8\)\,Hz while completing inference at
approximately \(0.74\)\,Hz, discarding roughly \(96\%\) of what it received.
Discarded frames are not free: each requires reassembly of some \(2{,}150\) UDP
fragments, deserialisation of \(2.76\)\,MB and a further \(2.76\)\,MB
colour-space conversion, all competing with inference for the same four
processor cores.

Reducing the camera rate at source from \(20\)\,Hz to \(5\)\,Hz --- still
approximately seven times the rate the detector can consume --- reduced inference
latency to approximately \(1{,}000\)\,ms, recovering about a quarter of the
deficit, and reduced link occupancy from \(481\)\,Mbps to \(111\)\,Mbps. The
application-level frame-rate control provided by the relay node cannot achieve
this, as it discards frames after reception, by which point the transport and
deserialisation cost has already been incurred.

This interaction is not observable in pure simulation, where sensor data is
delivered through shared memory at negligible cost, nor in an isolated benchmark
of the detector, since the interference originates outside it. It appears only
when a physical processor performs both reception and inference concurrently.

% ---------------------------------------------------------------------
\subsection{Limitations}
\label{sec:hitl-limitations}

Clock synchronisation between the two machines has not yet been established.
End-to-end latency is computed from timestamps taken on different hosts, and
without synchronisation those values are not meaningful while remaining plausible
in magnitude; synchronisation is required before end-to-end latency is reported.

The edge configuration is characterised; the corresponding ground configuration,
in which compressed imagery crosses the simulated channel and inference is
performed at the ground station, remains to be measured.

Measurements were taken with the aircraft close to the ground station, giving
zero packet loss throughout. Behaviour as range increases and buildings obstruct
the path has not been characterised. The testbed supports a single physical edge
node; multi-node operation has not been evaluated.
; cross-references
%  \ref{sec:hitl}. A longer version with the full 26-case test matrix is
%  kept in testing_validation_long.tex if an appendix is wanted.
%
%  REQUIRED PACKAGES
%      \usepackage{booktabs}
%      \usepackage{array}
%      \usepackage{graphicx}
%      \usepackage{subcaption}
%      \graphicspath{{figures/}}
%
%  IMAGES EXPECTED in docs/report/figures/  (see checklist at end)
%      detection_a.png  detection_b.png  detection_c.png
% =====================================================================

\section{Testing, Experimentation and Validation}
\label{sec:test}

Testing of the hardware-in-the-loop testbed described in
Section~\ref{sec:hitl} was organised around a difficulty specific to this
kind of system: a fault at a low layer does not raise an error at the layer
above it, but produces plausible and wrong behaviour there. A subscriber
starved of buffer space receives nothing and reports nothing; two nodes with
mismatched topic names run indefinitely without complaint. Testing was
therefore performed bottom-up --- physical link, path separation, middleware
transport, detector, end-to-end --- with each layer required to pass before
the next was attempted.

Three results are reported here. The first establishes that the two-link
architecture does what it claims. The second establishes that the detector
finds real people rather than plausible-looking texture. The third concerns
the measurement chain itself, and is included because it determines whether
the other two can be believed.

% ---------------------------------------------------------------------
\subsection{Test environment}
\label{sec:test-env}

Both machines run Ubuntu~22.04 with ROS\,2 Humble and Fast~DDS~2.6, so the
middleware is identical at each end of the link and cannot itself be a
source of difference. The simulation host runs Gazebo Classic in lockstep
with ArduPilot SITL, and ns-3 with TapBridge for the wireless channel. The
edge node is a Raspberry~Pi~4B running YOLOv8n restricted to the COCO
\emph{person} class. The camera produces $1280 \times 720$ frames of
$2.76$\,MB at $5$\,Hz. The test scene contains exactly five human models at
coordinates fixed in the world definition, which makes recall a counted
quantity rather than an estimated one.

The Pi's single Ethernet port is divided by 802.1Q tagging into
\texttt{eth0.10} ($10.0.0.2$, camera, unimpaired) and \texttt{eth0.42}
($10.42.0.12$, detections, through ns-3), created by
\texttt{scripts/netns/pi\_hitl\_link.sh}. The matching host interfaces
\texttt{eth-cam} and \texttt{eth-rf} are created by step~1c of
\texttt{scripts/netns/launch\_single\_uav\_netns.sh}, the latter bridged
into \texttt{tap-uav2}.

% ---------------------------------------------------------------------
\subsection{Result 1: the two links behave differently}
\label{sec:test-links}

The central claim of the architecture is that detection results cross a
simulated wireless channel while camera imagery does not. Two paths that
both work are not evidence of this; they must also be shown to differ in the
intended way, and imagery must be shown to be absent from the radio path.

Table~\ref{tab:test-links} reports sixty ICMP samples on the sensor path and
thirty on the radio path. Large samples are necessary: the first packet on the
radio path includes address resolution across the simulated channel, and a
two-sample measurement of the same configuration returned a mean of
$167$\,ms against $74.5$\,ms for thirty. The short sample is biased, not
merely noisy.

\begin{table}[htbp]
\centering
\caption{Measured characteristics of the two paths. Sensor path over sixty
samples, radio path over thirty; zero packet loss on both.}
\label{tab:test-links}
\begin{tabular}{@{}lrrrr@{}}
\toprule
Path & Min (ms) & Mean (ms) & Max (ms) & Jitter (ms) \\
\midrule
Sensor, VLAN 10 (wired)   & 1.0  & 1.4  & 1.8   & 0.12 \\
Radio, VLAN 42 (via ns-3) & 10.1 & 74.5 & 158.0 & 38.4 \\
\midrule
Ratio                     & --   & $53\times$ & -- & $320\times$ \\
\bottomrule
\end{tabular}
\end{table}

The radio path is approximately \(53\times\) slower on average, and its
variation is larger by more than two orders of magnitude.

Its internal spread is itself informative: a minimum of \(10.1\)\,ms against
a mean of \(74.5\)\,ms is a range of some \(7\times\) on one link, where a
channel applying a fixed delay would produce a narrow distribution. The
variation indicates that ns-3 is modelling contention, in which packets queue
and back off, rather than adding a constant offset.

The radio measurement was repeated. An independent set of thirty samples taken
several hours earlier, with the host reconfigured in between, gave a mean of
\(77.4\)\,ms against the \(74.5\)\,ms reported here, agreeing within
\(4\%\).

The sensor figures were taken \emph{while} the camera was streaming at
\(110\)\,Mbps. An idle-cable measurement over the same number of samples
gives the same mean, so the two logical links sharing one physical cable do
not delay one another. Queue management was verified rather than assumed:
both machines run \texttt{fq\_codel} with a \(5\)\,ms target, and byte
queue limits on the edge node had reduced its driver ring to approximately one
packet.

Absence of imagery on the radio path was established negatively, by
observation rather than by inspecting a rule set, since separation here is
enforced by addressing and bridge membership rather than by filtering. No
$10.0.0.0/24$ traffic was observed on \texttt{tap-uav2}, the point at which
traffic enters ns-3, and the byte counters on the two interfaces differ by
approximately four orders of magnitude --- consistent with $2.76$\,MB per
frame on one path against $118$\,B per detection message on the other, a
reduction of approximately \(23{,}000\times\).

A structural argument supports the measurements. The ground station receiver
executes inside the \texttt{gcsns} network namespace, whose only interface
is one end of a veth pair attached to \texttt{tap-gcs}. That namespace has
no interface on $10.0.0.0/24$ and no route to it. A message arriving at the
receiver cannot have taken any other path, because no other path exists from
inside the namespace. Separation is therefore a property of the topology
rather than of a setting that could be silently overridden.

% ---------------------------------------------------------------------
\subsection{Result 2: detections correspond to real subjects}
\label{sec:test-detection}

Visual confirmation was made a precondition of accepting any quantitative
result. A detector reporting a plausible count while placing its boxes on
building texture would satisfy every numerical check in this section, and an
earlier configuration did exactly that: sixteen confident detections against
five subjects, none correct. Only inspection of annotated frames
distinguishes the two cases. Those false positives were eliminated by
raising the confidence threshold and narrowing the camera field of view.

\begin{figure}[htbp]
\centering
\begin{subfigure}[b]{0.32\textwidth}
  \includegraphics[width=\textwidth]{detection_a.png}
  \caption{Two subjects, confidence $0.59$ and $0.42$.}
  \label{fig:test-det-a}
\end{subfigure}
\hfill
\begin{subfigure}[b]{0.32\textwidth}
  \includegraphics[width=\textwidth]{detection_b.png}
  \caption{One subject, confidence $0.48$.}
  \label{fig:test-det-b}
\end{subfigure}
\hfill
\begin{subfigure}[b]{0.32\textwidth}
  \includegraphics[width=\textwidth]{detection_c.png}
  \caption{Three subjects, confidence $0.59$, $0.47$ and $0.46$.}
  \label{fig:test-det-c}
\end{subfigure}
\caption{Detector output during flight, with the edge node's log above each
view. Bounding boxes coincide with the human models rather than with ground
texture. The log records every frame, including those with no detection, so
that a stalled node is distinguishable from an empty scene, and reports
inference duration per frame.}
\label{fig:test-detection}
\end{figure}

Figure~\ref{fig:test-detection} shows three moments from a patrol flight.
Between one and four of the five subjects were detected per frame, at
confidences of \(0.42\)--\(0.59\). The shortfall against the known count of
five is attributable to scene geometry rather than to detector deficiency:
the models stand approximately one metre apart, so from the oblique camera
view the nearer figures partially occlude those behind them, and a partially
occluded subject at approximately \(53\) pixels in height falls below the
confidence threshold. Because the number of subjects is known exactly, any
count exceeding five would be a false positive without the need to
adjudicate individual boxes; none was observed in this configuration.

Inference duration was measured around the model invocation alone, excluding
message reception, format conversion and publication, so that the figure
remains comparable between back-ends and against standalone execution. Over
\(480\) consecutive frames of one patrol flight the mean was
\(1{,}093\)\,ms, with a minimum of \(1{,}025\)\,ms. The distribution is tight
apart from a single outlier at \(6{,}949\)\,ms, one frame in \(480\),
attributable to contention for the processor from another task rather than to
the model itself.

The Pi therefore completes rather less than one detection per second against a
camera delivering five frames per second, and discards approximately \(80\%\)
of what it receives --- which is why the camera subscription uses
\texttt{BEST\_EFFORT} with queue depth one, keeping the newest frame rather
than accumulating a backlog of stale ones.

% ---------------------------------------------------------------------
\subsection{Result 3: validating the measurement chain}
\label{sec:test-chain}

The two results above depend on the camera actually reaching the edge node
at the rate the configuration specifies. That assumption failed, and the
manner of its failure is reported here because it determines what the other
measurements are worth.

The symptom was that the edge node received imagery at \(0.19\)\,Hz where
the camera was configured for \(5\)\,Hz, a shortfall of approximately
\(26\times\). The visible consequences appeared three layers above the
cause: the detector reported no humans in \(93\) of \(94\) frames, and the
patrol mission timed out at every waypoint. Both invite explanation in terms
of detector settings or flight control, and neither has anything to do with
either --- the aircraft flew correctly and the detector functioned
correctly, but the scene was being observed through roughly one frame every
five seconds.

\begin{table}[htbp]
\centering
\caption{Eliminating candidate causes of the frame-rate shortfall. Each row
is a measurement that excludes one layer.}
\label{tab:test-elimination}
\begin{tabular}{@{}p{38mm}p{38mm}p{40mm}@{}}
\toprule
Hypothesis & Measurement & Outcome \\
\midrule
Simulation running slowly
  & Real-time factor from \texttt{/clock}
  & \(0.998\); simulated time tracks wall clock \\[3pt]
Gazebo rendering slowly
  & Host-side subscriber on the camera topic
  & \(5.00\)\,Hz, gaps \(0.19\)--\(0.21\)\,s; source correct \\[3pt]
Link negotiated at $100$\,Mbps
  & \texttt{ethtool} on the edge node
  & $1000$\,Mb/s full duplex \\[3pt]
Fragments lost in flight
  & Interface counters, both machines
  & Zero errors, zero drops \\[3pt]
Receive buffers too small
  & \texttt{net.core.rmem\_max}
  & $512$\,MB, already raised \\[3pt]
\textbf{Send buffers too small}
  & \texttt{net.core.wmem\_max}
  & \(\mathbf{208}\)\,\textbf{kB}, left at the default \\
\bottomrule
\end{tabular}
\end{table}

The DDS profile requests a \(16\)\,MB send buffer. The kernel caps any such
request at \texttt{net.core.wmem\_max}, and performs the reduction silently.
A single frame is therefore approximately \(13\times\) larger than the
buffer available to transmit it: the publisher filled that buffer
immediately and discarded the remaining fragments at the socket layer,
before they reached the network interface. Comparing production with
transmission confirmed this --- the host generated \(414\)\,MB of imagery
over thirty seconds and placed \(10\)\,MB of it on the wire, \(2\)\,Mbps
against the \(110\)\,Mbps the configuration implies. Raising the parameter
on both machines and restarting the stack, which is required because
Fast~DDS applies socket options only at participant creation, restored
delivery to \(5.000\)\,Hz with an inter-arrival standard deviation of
\(4.5\)\,ms: an exact match to the source.

Two properties of this fault are worth recording. First, no counter
registers it. Loss at the socket layer is invisible to interface
statistics, and \texttt{tx\_dropped} on the sending interface and
\texttt{rx\_errors} on the receiving one both remained at zero throughout,
because the discarded fragments were never transmitted. A diagnosis founded
on interface counters would have concluded that the network was healthy,
which it was. Second, the two directions are governed by separate kernel
parameters. The receive path had been corrected during earlier work and the
transmit path had not, because the failure that prompted the original
correction pointed only at the receiving side. The pre-flight checks in
\texttt{scripts/netns/run\_hitl.sh} now verify both parameters on both
machines before the stack starts.

% ---------------------------------------------------------------------
\subsection{Limitations}
\label{sec:test-limitations}

Clock synchronisation between the two machines has not been established.
End-to-end latency would be computed from timestamps taken on different
hosts, and without synchronisation such values are not meaningful while
remaining plausible in magnitude. End-to-end latency is therefore not
reported, rather than reported with a caveat.

The edge configuration is characterised; the corresponding ground
configuration, in which compressed imagery crosses the simulated channel and
inference is performed at the ground station, remains to be measured. All
measurements were taken with the aircraft close to the ground station,
giving zero packet loss on both paths, so the reliability policy applied to
detection messages is untested under the loss conditions it exists to
handle. The testbed supports a single physical edge node; multi-node
operation has not been evaluated.

% =====================================================================
%  IMAGE CHECKLIST — save into docs/report/figures/
%
%  detection_a.png / detection_b.png / detection_c.png
%      The three screenshots already captured: the detector log above and
%      the rqt_image_view annotated frame below, showing 2, 1 and 3
%      detections respectively.
%
%      CROP EACH ONE before including. At three-to-a-row the full desktop
%      screenshot will be illegible. Keep only:
%        - the last ~8 log lines (enough to show the count and the
%          inference time), and
%        - the image pane itself, without the rqt toolbar and topic
%          selector.
%      Crop all three to the same aspect ratio so the row lines up.
%
%      If the log text is still too small at 0.32\textwidth, drop the log
%      from the figure entirely and show only the three annotated views —
%      the inference figures are already stated in the prose.
%
%  OPTIONAL, if a fourth image is wanted:
%      testbed_hardware.jpg — a photograph of the laptop, the USB gigabit
%      adapter, the single Ethernet cable and the Pi 4B. This is the only
%      image that shows the work is hardware-in-the-loop rather than pure
%      simulation. It would go in Section \ref{sec:test-env}.
% =====================================================================

%
%  Intended to follow % =====================================================================
%  Hardware-in-the-Loop Edge Processing
%  A SECTION of the FYP report. Include from the parent chapter with:
%      % =====================================================================
%  Hardware-in-the-Loop Edge Processing
%  A SECTION of the FYP report. Include from the parent chapter with:
%      \input{hitl_edge_processing}
%
%  HEADING LEVELS USED
%      \section     - the HITL section itself
%      \subsection  - its parts
%      \paragraph   - short findings within Results
%
%  REQUIRED PACKAGES (add to the main preamble if not already present)
%      \usepackage{booktabs}
%      \usepackage{graphicx}
%      \usepackage{tikz}
%      \usetikzlibrary{arrows.meta, positioning, fit, backgrounds}
%
%  All labels are prefixed hitl / tab:hitl- / fig:hitl- so they cannot
%  collide with labels elsewhere in the shared report.
% =====================================================================

\section{Hardware-in-the-Loop Edge Processing}
\label{sec:hitl}

The vision workload was moved off the simulation host and onto physical
hardware, so that the cost of onboard detection could be measured on a processor
representative of one that would actually fly. A Raspberry~Pi~4B performs
detection on imagery produced by the simulation, and returns its results across a
wireless channel simulated in ns-3.

% ---------------------------------------------------------------------
\subsection{Extending the testbed with physical hardware}
\label{sec:hitl-testbed}

The existing stack ran entirely on one machine: Gazebo for the environment and
sensors, ArduPilot SITL for the autopilot, ns-3 for the wireless channel, and
Linux network namespaces isolating the ground station from each aircraft. The
ns-3 scenario defined four nodes --- a ground control station and three aircraft
--- of which the second and third were unused placeholders.

A Raspberry~Pi~4B was introduced as the physical edge node and assigned to the
second aircraft slot. Its TAP device, previously created only to satisfy the
scenario's argument list, was brought into service as the attachment point. The
autopilot, physics and channel simulation remained on the host; only the vision
computation moved onto hardware.

The Pi was configured with Ubuntu~22.04 and ROS\,2 Humble to match the host
exactly, giving both machines the same DDS implementation. The host has no
Ethernet port, so a USB~3.0 gigabit adapter was fitted. An initial low-cost
adapter was replaced after \texttt{iperf3} measured \(4.5\)\,Mbps with \(24\%\)
loss, despite \texttt{ethtool} reporting \(100\)\,Mbps full duplex; the
replacement sustained \(834\)\,Mbps with no retransmissions.

% ---------------------------------------------------------------------
\subsection{Constructing the two-link network}
\label{sec:hitl-network}

Two logically separate paths were required between host and Pi. One carries
camera imagery and must remain unimpaired, since on a real aircraft it is a short
cable inside a single airframe rather than a radio. The other carries detection
results and must pass through the simulated channel, as that is the link the
experiment exists to characterise. The same reasoning applies to the
flight-dynamics traffic between SITL and Gazebo, which also bypasses ns-3.

The Pi provides a single Ethernet interface, so both paths were carried over one
cable using IEEE~802.1Q VLAN tagging. On the host, two virtual interfaces were
created on the physical adapter: \texttt{eth-cam}, tagged VLAN~10 and addressed
\(10.0.0.1\); and \texttt{eth-rf}, tagged VLAN~42, left without an address and
made a member of a new bridge, \texttt{br-uav2}. That bridge also received
\texttt{tap-uav2}, the ns-3 interface for the second aircraft node. On the Pi,
matching interfaces \texttt{eth0.10} (\(10.0.0.2\)) and \texttt{eth0.42}
(\(10.42.0.12\)) were created.

Attaching the physical interface to the bridge directly was rejected during
design. A Linux bridge floods frames to all member ports, so camera traffic would
have been delivered into the ns-3 TAP device as well, injecting approximately
\(111\)\,Mbps of imagery into the simulated channel.

The resulting arrangement enforces separation through routing rather than
filtering. Gazebo is addressed only on \(10.0.0.1\) and can reach the Pi only at
\(10.0.0.2\); the ground station namespace is addressed only on \(10.42.0.10\)
and can reach it only at \(10.42.0.12\). Neither can use the other's path.

\begin{figure}[htbp]
\centering
\begin{tikzpicture}[
  font=\small,
  box/.style={draw, rounded corners=2pt, align=center, inner sep=5pt},
  proc/.style={box, fill=black!4},
  net/.style={box, fill=black!10},
  arr/.style={-{Latex[length=2mm]}, thick},
]
\node[proc] (gz)   {Gazebo\\camera + physics};
\node[proc, below=14mm of gz] (gcs) {Ground station\\(namespace)\\\texttt{10.42.0.10}};
\node[net,  right=16mm of gcs] (ns3) {ns-3\\channel};
\node[net,  above=8mm of ns3]  (br)  {\texttt{br-uav2}\\+ TAP};
\node[proc, right=46mm of gz] (pi) {Raspberry Pi 4B\\YOLOv8n detector};

\draw[arr] (gz) -- node[above, align=center, font=\scriptsize]
  {VLAN 10 \quad camera, 2.76\,MB/frame\\unimpaired} (pi);
\draw[arr] (pi.south) |- node[pos=0.75, above, font=\scriptsize]
  {VLAN 42 \quad detections, 118\,B} (br.east);
\draw[arr] (br) -- (ns3);
\draw[arr] (ns3) -- (gcs);

\begin{scope}[on background layer]
  \node[draw, dashed, rounded corners, fit=(gz)(gcs)(ns3)(br),
        inner sep=7pt, label={[font=\scriptsize\itshape]above left:Host PC}] {};
\end{scope}
\end{tikzpicture}
\caption{The two-link arrangement. Camera imagery reaches the edge node over an
unimpaired wired path; detection results traverse the simulated channel. Both
share one physical cable, separated by VLAN tagging.}
\label{fig:hitl-architecture}
\end{figure}

Addressing was made persistent using NetworkManager on the host and netplan on
the Pi, the two distributions using different network managers. Addresses
assigned imperatively were found to be removed within minutes, as each manager
reconciles kernel state against its own configuration and deletes entries it did
not create.

% ---------------------------------------------------------------------
\subsection{Configuring the middleware}
\label{sec:hitl-dds}

A Fast DDS profile was written and applied on both machines. It pins participants
to the intended interfaces through an interface whitelist, since both machines
are attached to WiFi as well as the wired link and DDS otherwise advertises on
every available interface. The profile also disables the shared-memory transport,
which is not isolated by network namespaces and would allow participants to
exchange data without traversing the simulated channel.

Kernel socket buffers were raised on both machines from the \(208\)\,kB default
to \(512\)\,MB. A single \(1280 \times 720\) frame is \(2.76\)\,MB and fragments
into approximately \(1{,}900\) UDP datagrams, of which the default buffer holds
about \(7\%\); reassembly never completed and subscribers received nothing, with
no error reported.

The camera subscription on the Pi was changed from the ROS\,2 default of
\texttt{RELIABLE} with queue depth ten to \texttt{BEST\_EFFORT} with depth one.
Under the default policy the publisher retains each sample until the subscriber
acknowledges it, and because the Pi acknowledges at approximately \(1\)\,Hz, the
camera output of the entire simulation was throttled to \(0.27\)\,Hz.
\texttt{RELIABLE} was retained for detection messages, which must not be lost.

% ---------------------------------------------------------------------
\subsection{Deploying the detector}
\label{sec:hitl-detector}

The detection node runs YOLOv8n restricted to the COCO \emph{person} class,
subscribing to the camera topic and publishing results as JSON. It was deployed
to the Pi and instrumented to record inference duration, measured around the
model invocation alone so that figures remain comparable between back-ends and
against standalone execution.

Two dependency conflicts were resolved during deployment. Installing
\texttt{ultralytics} brings in NumPy~2.x, against which \texttt{cv\_bridge} --- a
compiled extension built for the 1.x C ABI --- fails. Installing an older OpenCV
to compensate reintroduces NumPy~2 as a dependency, and OpenCV~5.0 separately
breaks \texttt{cv\_bridge} by renaming the type constants it uses to build its
conversion table. Versions were pinned to \texttt{numpy==1.26.4} and
\texttt{opencv-python==4.10.0.84}, the only combination satisfying both.

The model was additionally exported to NCNN for ARM-optimised inference. Because
NCNN fixes the input tensor shape at export time, the export geometry was
specified explicitly as \(544 \times 960\) to match the letterboxed frame; a
default export produces a square input in which approximately \(43\%\) of the
computation is spent on padding, at \(19.9\)\,GFLOPS against \(11.3\)\,GFLOPS for
the correct shape. Half-precision export was avoided, the Cortex-A72 lacking fp16
SIMD support.

% ---------------------------------------------------------------------
\subsection{Automation}
\label{sec:hitl-automation}

Bring-up was consolidated into a single script that starts the namespaces, ns-3,
Gazebo, SITL, the micro-ROS agent and the ground station receiver, creates the
VLANs and bridge, and launches the detector on the Pi over SSH. Pre-flight checks
were added for each failure encountered during development: presence of the host
address before Gazebo starts, since Fast DDS enumerates interfaces only at
participant creation and cannot use one assigned later; reachability of the Pi;
adequate socket buffer limits on both machines; and confirmation after startup
that Gazebo bound the intended interface.

% ---------------------------------------------------------------------
\subsection{Results}
\label{sec:hitl-results}

Detections were verified visually, by inspecting annotated frames, before any
quantitative result was accepted. An earlier configuration had reported up to
sixteen detections against five subjects; these were false positives on building
texture, eliminated by raising the confidence threshold and narrowing the camera
field of view.

\paragraph{Channel characterisation.}
Table~\ref{tab:hitl-links} compares the two paths, each measured over thirty ICMP
samples. Short samples are unreliable: the first packet includes address
resolution across the simulated channel, and a two-sample measurement of the same
configuration returned a mean of \(167\)\,ms against \(59\)\,ms for thirty.

\begin{table}[htbp]
\centering
\caption{Measured characteristics of the two links, thirty samples each.}
\label{tab:hitl-links}
\begin{tabular}{lrrrr}
\toprule
Path & Min (ms) & Mean (ms) & Max (ms) & Jitter (ms) \\
\midrule
Sensor, VLAN 10 (wired)   & 0.4  & 1.5  & 2.7   & 0.4  \\
Radio, VLAN 42 (via ns-3) & 11.7 & 77.4 & 242.6 & 54.8 \\
\bottomrule
\end{tabular}
\end{table}

Packet loss was zero on both paths, the aircraft remaining close to the ground
station throughout. The radio path is approximately \(53\times\) slower than the
wired path on average and varies approximately \(140\times\) more.

The distribution within the radio path is itself informative. Its minimum of
\(11.7\)\,ms against a mean of \(77.4\)\,ms is a spread of some \(6.6\times\) on
one link, where a channel applying a fixed delay would produce a narrow one.
The variation indicates that ns-3 is modelling contention, in which packets
queue and back off before transmission, rather than simply adding a constant
offset to each packet.

\paragraph{Sensor transport.}
With the camera at \(20\)\,Hz the wired path carried \(481\)\,Mbps at
\(42{,}544\) packets per second with zero dropped packets, occupying
approximately \(58\%\) of measured link capacity.

\paragraph{Inference latency.}
Table~\ref{tab:hitl-inference} reports latency for each combination of back-end
and input resolution. Changing back-end at constant resolution gives
approximately \(2.5\times\); halving the pixel count at constant back-end gives
approximately \(2.3\times\). Latency therefore scales close to linearly with
pixel count, a pixel ratio of \(0.47\) producing a time ratio of \(0.43\).

\begin{table}[htbp]
\centering
\caption{Inference latency on the Raspberry Pi 4B.}
\label{tab:hitl-inference}
\begin{tabular}{llrr}
\toprule
Back-end & Input & Latency (ms) & Rate (fps) \\
\midrule
PyTorch & $960 \times 544$ & 2540 & 0.39 \\
PyTorch & $640 \times 384$ & 1100 & 0.91 \\
NCNN    & $960 \times 544$ & 1000 & 1.00 \\
NCNN    & $640 \times 384$ & \multicolumn{2}{l}{not measured (est.\ 470)} \\
\bottomrule
\end{tabular}
\end{table}

\paragraph{Detection performance.}
Between two and four of the five subjects were detected per frame, a recall of
\(40\)--\(80\%\). The shortfall is attributable to scene geometry rather than
detector deficiency: the models are separated by approximately one metre, so from
the oblique camera view the nearer figures partially occlude those behind them,
and partially occluded subjects at approximately \(53\) pixels in height fall
below the confidence threshold. Because the number and position of subjects are
known exactly from the world definition, recall is a directly measurable quantity
rather than an estimate.

\paragraph{Transmitted volume.}
Each detection message is approximately \(118\)\,B, against \(2.76\)\,MB for the
corresponding frame, a reduction of approximately \(23{,}000\times\) in data
crossing the radio.

% ---------------------------------------------------------------------
\subsection{Contention between sensor transport and inference}
\label{sec:hitl-contention}

The two workloads on the edge node proved not to be independent. The same model,
on the same input, required \(623\)\,ms executed on a static image but
\(1{,}343\)\,ms within the live pipeline: a difference of \(720\)\,ms, or
approximately half the available inference throughput.

The edge node received frames at \(19.8\)\,Hz while completing inference at
approximately \(0.74\)\,Hz, discarding roughly \(96\%\) of what it received.
Discarded frames are not free: each requires reassembly of some \(2{,}150\) UDP
fragments, deserialisation of \(2.76\)\,MB and a further \(2.76\)\,MB
colour-space conversion, all competing with inference for the same four
processor cores.

Reducing the camera rate at source from \(20\)\,Hz to \(5\)\,Hz --- still
approximately seven times the rate the detector can consume --- reduced inference
latency to approximately \(1{,}000\)\,ms, recovering about a quarter of the
deficit, and reduced link occupancy from \(481\)\,Mbps to \(111\)\,Mbps. The
application-level frame-rate control provided by the relay node cannot achieve
this, as it discards frames after reception, by which point the transport and
deserialisation cost has already been incurred.

This interaction is not observable in pure simulation, where sensor data is
delivered through shared memory at negligible cost, nor in an isolated benchmark
of the detector, since the interference originates outside it. It appears only
when a physical processor performs both reception and inference concurrently.

% ---------------------------------------------------------------------
\subsection{Limitations}
\label{sec:hitl-limitations}

Clock synchronisation between the two machines has not yet been established.
End-to-end latency is computed from timestamps taken on different hosts, and
without synchronisation those values are not meaningful while remaining plausible
in magnitude; synchronisation is required before end-to-end latency is reported.

The edge configuration is characterised; the corresponding ground configuration,
in which compressed imagery crosses the simulated channel and inference is
performed at the ground station, remains to be measured.

Measurements were taken with the aircraft close to the ground station, giving
zero packet loss throughout. Behaviour as range increases and buildings obstruct
the path has not been characterised. The testbed supports a single physical edge
node; multi-node operation has not been evaluated.

%
%  HEADING LEVELS USED
%      \section     - the HITL section itself
%      \subsection  - its parts
%      \paragraph   - short findings within Results
%
%  REQUIRED PACKAGES (add to the main preamble if not already present)
%      \usepackage{booktabs}
%      \usepackage{graphicx}
%      \usepackage{tikz}
%      \usetikzlibrary{arrows.meta, positioning, fit, backgrounds}
%
%  All labels are prefixed hitl / tab:hitl- / fig:hitl- so they cannot
%  collide with labels elsewhere in the shared report.
% =====================================================================

\section{Hardware-in-the-Loop Edge Processing}
\label{sec:hitl}

The vision workload was moved off the simulation host and onto physical
hardware, so that the cost of onboard detection could be measured on a processor
representative of one that would actually fly. A Raspberry~Pi~4B performs
detection on imagery produced by the simulation, and returns its results across a
wireless channel simulated in ns-3.

% ---------------------------------------------------------------------
\subsection{Extending the testbed with physical hardware}
\label{sec:hitl-testbed}

The existing stack ran entirely on one machine: Gazebo for the environment and
sensors, ArduPilot SITL for the autopilot, ns-3 for the wireless channel, and
Linux network namespaces isolating the ground station from each aircraft. The
ns-3 scenario defined four nodes --- a ground control station and three aircraft
--- of which the second and third were unused placeholders.

A Raspberry~Pi~4B was introduced as the physical edge node and assigned to the
second aircraft slot. Its TAP device, previously created only to satisfy the
scenario's argument list, was brought into service as the attachment point. The
autopilot, physics and channel simulation remained on the host; only the vision
computation moved onto hardware.

The Pi was configured with Ubuntu~22.04 and ROS\,2 Humble to match the host
exactly, giving both machines the same DDS implementation. The host has no
Ethernet port, so a USB~3.0 gigabit adapter was fitted. An initial low-cost
adapter was replaced after \texttt{iperf3} measured \(4.5\)\,Mbps with \(24\%\)
loss, despite \texttt{ethtool} reporting \(100\)\,Mbps full duplex; the
replacement sustained \(834\)\,Mbps with no retransmissions.

% ---------------------------------------------------------------------
\subsection{Constructing the two-link network}
\label{sec:hitl-network}

Two logically separate paths were required between host and Pi. One carries
camera imagery and must remain unimpaired, since on a real aircraft it is a short
cable inside a single airframe rather than a radio. The other carries detection
results and must pass through the simulated channel, as that is the link the
experiment exists to characterise. The same reasoning applies to the
flight-dynamics traffic between SITL and Gazebo, which also bypasses ns-3.

The Pi provides a single Ethernet interface, so both paths were carried over one
cable using IEEE~802.1Q VLAN tagging. On the host, two virtual interfaces were
created on the physical adapter: \texttt{eth-cam}, tagged VLAN~10 and addressed
\(10.0.0.1\); and \texttt{eth-rf}, tagged VLAN~42, left without an address and
made a member of a new bridge, \texttt{br-uav2}. That bridge also received
\texttt{tap-uav2}, the ns-3 interface for the second aircraft node. On the Pi,
matching interfaces \texttt{eth0.10} (\(10.0.0.2\)) and \texttt{eth0.42}
(\(10.42.0.12\)) were created.

Attaching the physical interface to the bridge directly was rejected during
design. A Linux bridge floods frames to all member ports, so camera traffic would
have been delivered into the ns-3 TAP device as well, injecting approximately
\(111\)\,Mbps of imagery into the simulated channel.

The resulting arrangement enforces separation through routing rather than
filtering. Gazebo is addressed only on \(10.0.0.1\) and can reach the Pi only at
\(10.0.0.2\); the ground station namespace is addressed only on \(10.42.0.10\)
and can reach it only at \(10.42.0.12\). Neither can use the other's path.

\begin{figure}[htbp]
\centering
\begin{tikzpicture}[
  font=\small,
  box/.style={draw, rounded corners=2pt, align=center, inner sep=5pt},
  proc/.style={box, fill=black!4},
  net/.style={box, fill=black!10},
  arr/.style={-{Latex[length=2mm]}, thick},
]
\node[proc] (gz)   {Gazebo\\camera + physics};
\node[proc, below=14mm of gz] (gcs) {Ground station\\(namespace)\\\texttt{10.42.0.10}};
\node[net,  right=16mm of gcs] (ns3) {ns-3\\channel};
\node[net,  above=8mm of ns3]  (br)  {\texttt{br-uav2}\\+ TAP};
\node[proc, right=46mm of gz] (pi) {Raspberry Pi 4B\\YOLOv8n detector};

\draw[arr] (gz) -- node[above, align=center, font=\scriptsize]
  {VLAN 10 \quad camera, 2.76\,MB/frame\\unimpaired} (pi);
\draw[arr] (pi.south) |- node[pos=0.75, above, font=\scriptsize]
  {VLAN 42 \quad detections, 118\,B} (br.east);
\draw[arr] (br) -- (ns3);
\draw[arr] (ns3) -- (gcs);

\begin{scope}[on background layer]
  \node[draw, dashed, rounded corners, fit=(gz)(gcs)(ns3)(br),
        inner sep=7pt, label={[font=\scriptsize\itshape]above left:Host PC}] {};
\end{scope}
\end{tikzpicture}
\caption{The two-link arrangement. Camera imagery reaches the edge node over an
unimpaired wired path; detection results traverse the simulated channel. Both
share one physical cable, separated by VLAN tagging.}
\label{fig:hitl-architecture}
\end{figure}

Addressing was made persistent using NetworkManager on the host and netplan on
the Pi, the two distributions using different network managers. Addresses
assigned imperatively were found to be removed within minutes, as each manager
reconciles kernel state against its own configuration and deletes entries it did
not create.

% ---------------------------------------------------------------------
\subsection{Configuring the middleware}
\label{sec:hitl-dds}

A Fast DDS profile was written and applied on both machines. It pins participants
to the intended interfaces through an interface whitelist, since both machines
are attached to WiFi as well as the wired link and DDS otherwise advertises on
every available interface. The profile also disables the shared-memory transport,
which is not isolated by network namespaces and would allow participants to
exchange data without traversing the simulated channel.

Kernel socket buffers were raised on both machines from the \(208\)\,kB default
to \(512\)\,MB. A single \(1280 \times 720\) frame is \(2.76\)\,MB and fragments
into approximately \(1{,}900\) UDP datagrams, of which the default buffer holds
about \(7\%\); reassembly never completed and subscribers received nothing, with
no error reported.

The camera subscription on the Pi was changed from the ROS\,2 default of
\texttt{RELIABLE} with queue depth ten to \texttt{BEST\_EFFORT} with depth one.
Under the default policy the publisher retains each sample until the subscriber
acknowledges it, and because the Pi acknowledges at approximately \(1\)\,Hz, the
camera output of the entire simulation was throttled to \(0.27\)\,Hz.
\texttt{RELIABLE} was retained for detection messages, which must not be lost.

% ---------------------------------------------------------------------
\subsection{Deploying the detector}
\label{sec:hitl-detector}

The detection node runs YOLOv8n restricted to the COCO \emph{person} class,
subscribing to the camera topic and publishing results as JSON. It was deployed
to the Pi and instrumented to record inference duration, measured around the
model invocation alone so that figures remain comparable between back-ends and
against standalone execution.

Two dependency conflicts were resolved during deployment. Installing
\texttt{ultralytics} brings in NumPy~2.x, against which \texttt{cv\_bridge} --- a
compiled extension built for the 1.x C ABI --- fails. Installing an older OpenCV
to compensate reintroduces NumPy~2 as a dependency, and OpenCV~5.0 separately
breaks \texttt{cv\_bridge} by renaming the type constants it uses to build its
conversion table. Versions were pinned to \texttt{numpy==1.26.4} and
\texttt{opencv-python==4.10.0.84}, the only combination satisfying both.

The model was additionally exported to NCNN for ARM-optimised inference. Because
NCNN fixes the input tensor shape at export time, the export geometry was
specified explicitly as \(544 \times 960\) to match the letterboxed frame; a
default export produces a square input in which approximately \(43\%\) of the
computation is spent on padding, at \(19.9\)\,GFLOPS against \(11.3\)\,GFLOPS for
the correct shape. Half-precision export was avoided, the Cortex-A72 lacking fp16
SIMD support.

% ---------------------------------------------------------------------
\subsection{Automation}
\label{sec:hitl-automation}

Bring-up was consolidated into a single script that starts the namespaces, ns-3,
Gazebo, SITL, the micro-ROS agent and the ground station receiver, creates the
VLANs and bridge, and launches the detector on the Pi over SSH. Pre-flight checks
were added for each failure encountered during development: presence of the host
address before Gazebo starts, since Fast DDS enumerates interfaces only at
participant creation and cannot use one assigned later; reachability of the Pi;
adequate socket buffer limits on both machines; and confirmation after startup
that Gazebo bound the intended interface.

% ---------------------------------------------------------------------
\subsection{Results}
\label{sec:hitl-results}

Detections were verified visually, by inspecting annotated frames, before any
quantitative result was accepted. An earlier configuration had reported up to
sixteen detections against five subjects; these were false positives on building
texture, eliminated by raising the confidence threshold and narrowing the camera
field of view.

\paragraph{Channel characterisation.}
Table~\ref{tab:hitl-links} compares the two paths, each measured over thirty ICMP
samples. Short samples are unreliable: the first packet includes address
resolution across the simulated channel, and a two-sample measurement of the same
configuration returned a mean of \(167\)\,ms against \(59\)\,ms for thirty.

\begin{table}[htbp]
\centering
\caption{Measured characteristics of the two links, thirty samples each.}
\label{tab:hitl-links}
\begin{tabular}{lrrrr}
\toprule
Path & Min (ms) & Mean (ms) & Max (ms) & Jitter (ms) \\
\midrule
Sensor, VLAN 10 (wired)   & 0.4  & 1.5  & 2.7   & 0.4  \\
Radio, VLAN 42 (via ns-3) & 11.7 & 77.4 & 242.6 & 54.8 \\
\bottomrule
\end{tabular}
\end{table}

Packet loss was zero on both paths, the aircraft remaining close to the ground
station throughout. The radio path is approximately \(53\times\) slower than the
wired path on average and varies approximately \(140\times\) more.

The distribution within the radio path is itself informative. Its minimum of
\(11.7\)\,ms against a mean of \(77.4\)\,ms is a spread of some \(6.6\times\) on
one link, where a channel applying a fixed delay would produce a narrow one.
The variation indicates that ns-3 is modelling contention, in which packets
queue and back off before transmission, rather than simply adding a constant
offset to each packet.

\paragraph{Sensor transport.}
With the camera at \(20\)\,Hz the wired path carried \(481\)\,Mbps at
\(42{,}544\) packets per second with zero dropped packets, occupying
approximately \(58\%\) of measured link capacity.

\paragraph{Inference latency.}
Table~\ref{tab:hitl-inference} reports latency for each combination of back-end
and input resolution. Changing back-end at constant resolution gives
approximately \(2.5\times\); halving the pixel count at constant back-end gives
approximately \(2.3\times\). Latency therefore scales close to linearly with
pixel count, a pixel ratio of \(0.47\) producing a time ratio of \(0.43\).

\begin{table}[htbp]
\centering
\caption{Inference latency on the Raspberry Pi 4B.}
\label{tab:hitl-inference}
\begin{tabular}{llrr}
\toprule
Back-end & Input & Latency (ms) & Rate (fps) \\
\midrule
PyTorch & $960 \times 544$ & 2540 & 0.39 \\
PyTorch & $640 \times 384$ & 1100 & 0.91 \\
NCNN    & $960 \times 544$ & 1000 & 1.00 \\
NCNN    & $640 \times 384$ & \multicolumn{2}{l}{not measured (est.\ 470)} \\
\bottomrule
\end{tabular}
\end{table}

\paragraph{Detection performance.}
Between two and four of the five subjects were detected per frame, a recall of
\(40\)--\(80\%\). The shortfall is attributable to scene geometry rather than
detector deficiency: the models are separated by approximately one metre, so from
the oblique camera view the nearer figures partially occlude those behind them,
and partially occluded subjects at approximately \(53\) pixels in height fall
below the confidence threshold. Because the number and position of subjects are
known exactly from the world definition, recall is a directly measurable quantity
rather than an estimate.

\paragraph{Transmitted volume.}
Each detection message is approximately \(118\)\,B, against \(2.76\)\,MB for the
corresponding frame, a reduction of approximately \(23{,}000\times\) in data
crossing the radio.

% ---------------------------------------------------------------------
\subsection{Contention between sensor transport and inference}
\label{sec:hitl-contention}

The two workloads on the edge node proved not to be independent. The same model,
on the same input, required \(623\)\,ms executed on a static image but
\(1{,}343\)\,ms within the live pipeline: a difference of \(720\)\,ms, or
approximately half the available inference throughput.

The edge node received frames at \(19.8\)\,Hz while completing inference at
approximately \(0.74\)\,Hz, discarding roughly \(96\%\) of what it received.
Discarded frames are not free: each requires reassembly of some \(2{,}150\) UDP
fragments, deserialisation of \(2.76\)\,MB and a further \(2.76\)\,MB
colour-space conversion, all competing with inference for the same four
processor cores.

Reducing the camera rate at source from \(20\)\,Hz to \(5\)\,Hz --- still
approximately seven times the rate the detector can consume --- reduced inference
latency to approximately \(1{,}000\)\,ms, recovering about a quarter of the
deficit, and reduced link occupancy from \(481\)\,Mbps to \(111\)\,Mbps. The
application-level frame-rate control provided by the relay node cannot achieve
this, as it discards frames after reception, by which point the transport and
deserialisation cost has already been incurred.

This interaction is not observable in pure simulation, where sensor data is
delivered through shared memory at negligible cost, nor in an isolated benchmark
of the detector, since the interference originates outside it. It appears only
when a physical processor performs both reception and inference concurrently.

% ---------------------------------------------------------------------
\subsection{Limitations}
\label{sec:hitl-limitations}

Clock synchronisation between the two machines has not yet been established.
End-to-end latency is computed from timestamps taken on different hosts, and
without synchronisation those values are not meaningful while remaining plausible
in magnitude; synchronisation is required before end-to-end latency is reported.

The edge configuration is characterised; the corresponding ground configuration,
in which compressed imagery crosses the simulated channel and inference is
performed at the ground station, remains to be measured.

Measurements were taken with the aircraft close to the ground station, giving
zero packet loss throughout. Behaviour as range increases and buildings obstruct
the path has not been characterised. The testbed supports a single physical edge
node; multi-node operation has not been evaluated.
; cross-references
%  \ref{sec:hitl}. A longer version with the full 26-case test matrix is
%  kept in testing_validation_long.tex if an appendix is wanted.
%
%  REQUIRED PACKAGES
%      \usepackage{booktabs}
%      \usepackage{array}
%      \usepackage{graphicx}
%      \usepackage{subcaption}
%      \graphicspath{{figures/}}
%
%  IMAGES EXPECTED in docs/report/figures/  (see checklist at end)
%      detection_a.png  detection_b.png  detection_c.png
% =====================================================================

\section{Testing, Experimentation and Validation}
\label{sec:test}

Testing of the hardware-in-the-loop testbed described in
Section~\ref{sec:hitl} was organised around a difficulty specific to this
kind of system: a fault at a low layer does not raise an error at the layer
above it, but produces plausible and wrong behaviour there. A subscriber
starved of buffer space receives nothing and reports nothing; two nodes with
mismatched topic names run indefinitely without complaint. Testing was
therefore performed bottom-up --- physical link, path separation, middleware
transport, detector, end-to-end --- with each layer required to pass before
the next was attempted.

Three results are reported here. The first establishes that the two-link
architecture does what it claims. The second establishes that the detector
finds real people rather than plausible-looking texture. The third concerns
the measurement chain itself, and is included because it determines whether
the other two can be believed.

% ---------------------------------------------------------------------
\subsection{Test environment}
\label{sec:test-env}

Both machines run Ubuntu~22.04 with ROS\,2 Humble and Fast~DDS~2.6, so the
middleware is identical at each end of the link and cannot itself be a
source of difference. The simulation host runs Gazebo Classic in lockstep
with ArduPilot SITL, and ns-3 with TapBridge for the wireless channel. The
edge node is a Raspberry~Pi~4B running YOLOv8n restricted to the COCO
\emph{person} class. The camera produces $1280 \times 720$ frames of
$2.76$\,MB at $5$\,Hz. The test scene contains exactly five human models at
coordinates fixed in the world definition, which makes recall a counted
quantity rather than an estimated one.

The Pi's single Ethernet port is divided by 802.1Q tagging into
\texttt{eth0.10} ($10.0.0.2$, camera, unimpaired) and \texttt{eth0.42}
($10.42.0.12$, detections, through ns-3), created by
\texttt{scripts/netns/pi\_hitl\_link.sh}. The matching host interfaces
\texttt{eth-cam} and \texttt{eth-rf} are created by step~1c of
\texttt{scripts/netns/launch\_single\_uav\_netns.sh}, the latter bridged
into \texttt{tap-uav2}.

% ---------------------------------------------------------------------
\subsection{Result 1: the two links behave differently}
\label{sec:test-links}

The central claim of the architecture is that detection results cross a
simulated wireless channel while camera imagery does not. Two paths that
both work are not evidence of this; they must also be shown to differ in the
intended way, and imagery must be shown to be absent from the radio path.

Table~\ref{tab:test-links} reports sixty ICMP samples on the sensor path and
thirty on the radio path. Large samples are necessary: the first packet on the
radio path includes address resolution across the simulated channel, and a
two-sample measurement of the same configuration returned a mean of
$167$\,ms against $74.5$\,ms for thirty. The short sample is biased, not
merely noisy.

\begin{table}[htbp]
\centering
\caption{Measured characteristics of the two paths. Sensor path over sixty
samples, radio path over thirty; zero packet loss on both.}
\label{tab:test-links}
\begin{tabular}{@{}lrrrr@{}}
\toprule
Path & Min (ms) & Mean (ms) & Max (ms) & Jitter (ms) \\
\midrule
Sensor, VLAN 10 (wired)   & 1.0  & 1.4  & 1.8   & 0.12 \\
Radio, VLAN 42 (via ns-3) & 10.1 & 74.5 & 158.0 & 38.4 \\
\midrule
Ratio                     & --   & $53\times$ & -- & $320\times$ \\
\bottomrule
\end{tabular}
\end{table}

The radio path is approximately \(53\times\) slower on average, and its
variation is larger by more than two orders of magnitude.

Its internal spread is itself informative: a minimum of \(10.1\)\,ms against
a mean of \(74.5\)\,ms is a range of some \(7\times\) on one link, where a
channel applying a fixed delay would produce a narrow distribution. The
variation indicates that ns-3 is modelling contention, in which packets queue
and back off, rather than adding a constant offset.

The radio measurement was repeated. An independent set of thirty samples taken
several hours earlier, with the host reconfigured in between, gave a mean of
\(77.4\)\,ms against the \(74.5\)\,ms reported here, agreeing within
\(4\%\).

The sensor figures were taken \emph{while} the camera was streaming at
\(110\)\,Mbps. An idle-cable measurement over the same number of samples
gives the same mean, so the two logical links sharing one physical cable do
not delay one another. Queue management was verified rather than assumed:
both machines run \texttt{fq\_codel} with a \(5\)\,ms target, and byte
queue limits on the edge node had reduced its driver ring to approximately one
packet.

Absence of imagery on the radio path was established negatively, by
observation rather than by inspecting a rule set, since separation here is
enforced by addressing and bridge membership rather than by filtering. No
$10.0.0.0/24$ traffic was observed on \texttt{tap-uav2}, the point at which
traffic enters ns-3, and the byte counters on the two interfaces differ by
approximately four orders of magnitude --- consistent with $2.76$\,MB per
frame on one path against $118$\,B per detection message on the other, a
reduction of approximately \(23{,}000\times\).

A structural argument supports the measurements. The ground station receiver
executes inside the \texttt{gcsns} network namespace, whose only interface
is one end of a veth pair attached to \texttt{tap-gcs}. That namespace has
no interface on $10.0.0.0/24$ and no route to it. A message arriving at the
receiver cannot have taken any other path, because no other path exists from
inside the namespace. Separation is therefore a property of the topology
rather than of a setting that could be silently overridden.

% ---------------------------------------------------------------------
\subsection{Result 2: detections correspond to real subjects}
\label{sec:test-detection}

Visual confirmation was made a precondition of accepting any quantitative
result. A detector reporting a plausible count while placing its boxes on
building texture would satisfy every numerical check in this section, and an
earlier configuration did exactly that: sixteen confident detections against
five subjects, none correct. Only inspection of annotated frames
distinguishes the two cases. Those false positives were eliminated by
raising the confidence threshold and narrowing the camera field of view.

\begin{figure}[htbp]
\centering
\begin{subfigure}[b]{0.32\textwidth}
  \includegraphics[width=\textwidth]{detection_a.png}
  \caption{Two subjects, confidence $0.59$ and $0.42$.}
  \label{fig:test-det-a}
\end{subfigure}
\hfill
\begin{subfigure}[b]{0.32\textwidth}
  \includegraphics[width=\textwidth]{detection_b.png}
  \caption{One subject, confidence $0.48$.}
  \label{fig:test-det-b}
\end{subfigure}
\hfill
\begin{subfigure}[b]{0.32\textwidth}
  \includegraphics[width=\textwidth]{detection_c.png}
  \caption{Three subjects, confidence $0.59$, $0.47$ and $0.46$.}
  \label{fig:test-det-c}
\end{subfigure}
\caption{Detector output during flight, with the edge node's log above each
view. Bounding boxes coincide with the human models rather than with ground
texture. The log records every frame, including those with no detection, so
that a stalled node is distinguishable from an empty scene, and reports
inference duration per frame.}
\label{fig:test-detection}
\end{figure}

Figure~\ref{fig:test-detection} shows three moments from a patrol flight.
Between one and four of the five subjects were detected per frame, at
confidences of \(0.42\)--\(0.59\). The shortfall against the known count of
five is attributable to scene geometry rather than to detector deficiency:
the models stand approximately one metre apart, so from the oblique camera
view the nearer figures partially occlude those behind them, and a partially
occluded subject at approximately \(53\) pixels in height falls below the
confidence threshold. Because the number of subjects is known exactly, any
count exceeding five would be a false positive without the need to
adjudicate individual boxes; none was observed in this configuration.

Inference duration was measured around the model invocation alone, excluding
message reception, format conversion and publication, so that the figure
remains comparable between back-ends and against standalone execution. Over
\(480\) consecutive frames of one patrol flight the mean was
\(1{,}093\)\,ms, with a minimum of \(1{,}025\)\,ms. The distribution is tight
apart from a single outlier at \(6{,}949\)\,ms, one frame in \(480\),
attributable to contention for the processor from another task rather than to
the model itself.

The Pi therefore completes rather less than one detection per second against a
camera delivering five frames per second, and discards approximately \(80\%\)
of what it receives --- which is why the camera subscription uses
\texttt{BEST\_EFFORT} with queue depth one, keeping the newest frame rather
than accumulating a backlog of stale ones.

% ---------------------------------------------------------------------
\subsection{Result 3: validating the measurement chain}
\label{sec:test-chain}

The two results above depend on the camera actually reaching the edge node
at the rate the configuration specifies. That assumption failed, and the
manner of its failure is reported here because it determines what the other
measurements are worth.

The symptom was that the edge node received imagery at \(0.19\)\,Hz where
the camera was configured for \(5\)\,Hz, a shortfall of approximately
\(26\times\). The visible consequences appeared three layers above the
cause: the detector reported no humans in \(93\) of \(94\) frames, and the
patrol mission timed out at every waypoint. Both invite explanation in terms
of detector settings or flight control, and neither has anything to do with
either --- the aircraft flew correctly and the detector functioned
correctly, but the scene was being observed through roughly one frame every
five seconds.

\begin{table}[htbp]
\centering
\caption{Eliminating candidate causes of the frame-rate shortfall. Each row
is a measurement that excludes one layer.}
\label{tab:test-elimination}
\begin{tabular}{@{}p{38mm}p{38mm}p{40mm}@{}}
\toprule
Hypothesis & Measurement & Outcome \\
\midrule
Simulation running slowly
  & Real-time factor from \texttt{/clock}
  & \(0.998\); simulated time tracks wall clock \\[3pt]
Gazebo rendering slowly
  & Host-side subscriber on the camera topic
  & \(5.00\)\,Hz, gaps \(0.19\)--\(0.21\)\,s; source correct \\[3pt]
Link negotiated at $100$\,Mbps
  & \texttt{ethtool} on the edge node
  & $1000$\,Mb/s full duplex \\[3pt]
Fragments lost in flight
  & Interface counters, both machines
  & Zero errors, zero drops \\[3pt]
Receive buffers too small
  & \texttt{net.core.rmem\_max}
  & $512$\,MB, already raised \\[3pt]
\textbf{Send buffers too small}
  & \texttt{net.core.wmem\_max}
  & \(\mathbf{208}\)\,\textbf{kB}, left at the default \\
\bottomrule
\end{tabular}
\end{table}

The DDS profile requests a \(16\)\,MB send buffer. The kernel caps any such
request at \texttt{net.core.wmem\_max}, and performs the reduction silently.
A single frame is therefore approximately \(13\times\) larger than the
buffer available to transmit it: the publisher filled that buffer
immediately and discarded the remaining fragments at the socket layer,
before they reached the network interface. Comparing production with
transmission confirmed this --- the host generated \(414\)\,MB of imagery
over thirty seconds and placed \(10\)\,MB of it on the wire, \(2\)\,Mbps
against the \(110\)\,Mbps the configuration implies. Raising the parameter
on both machines and restarting the stack, which is required because
Fast~DDS applies socket options only at participant creation, restored
delivery to \(5.000\)\,Hz with an inter-arrival standard deviation of
\(4.5\)\,ms: an exact match to the source.

Two properties of this fault are worth recording. First, no counter
registers it. Loss at the socket layer is invisible to interface
statistics, and \texttt{tx\_dropped} on the sending interface and
\texttt{rx\_errors} on the receiving one both remained at zero throughout,
because the discarded fragments were never transmitted. A diagnosis founded
on interface counters would have concluded that the network was healthy,
which it was. Second, the two directions are governed by separate kernel
parameters. The receive path had been corrected during earlier work and the
transmit path had not, because the failure that prompted the original
correction pointed only at the receiving side. The pre-flight checks in
\texttt{scripts/netns/run\_hitl.sh} now verify both parameters on both
machines before the stack starts.

% ---------------------------------------------------------------------
\subsection{Limitations}
\label{sec:test-limitations}

Clock synchronisation between the two machines has not been established.
End-to-end latency would be computed from timestamps taken on different
hosts, and without synchronisation such values are not meaningful while
remaining plausible in magnitude. End-to-end latency is therefore not
reported, rather than reported with a caveat.

The edge configuration is characterised; the corresponding ground
configuration, in which compressed imagery crosses the simulated channel and
inference is performed at the ground station, remains to be measured. All
measurements were taken with the aircraft close to the ground station,
giving zero packet loss on both paths, so the reliability policy applied to
detection messages is untested under the loss conditions it exists to
handle. The testbed supports a single physical edge node; multi-node
operation has not been evaluated.

% =====================================================================
%  IMAGE CHECKLIST — save into docs/report/figures/
%
%  detection_a.png / detection_b.png / detection_c.png
%      The three screenshots already captured: the detector log above and
%      the rqt_image_view annotated frame below, showing 2, 1 and 3
%      detections respectively.
%
%      CROP EACH ONE before including. At three-to-a-row the full desktop
%      screenshot will be illegible. Keep only:
%        - the last ~8 log lines (enough to show the count and the
%          inference time), and
%        - the image pane itself, without the rqt toolbar and topic
%          selector.
%      Crop all three to the same aspect ratio so the row lines up.
%
%      If the log text is still too small at 0.32\textwidth, drop the log
%      from the figure entirely and show only the three annotated views —
%      the inference figures are already stated in the prose.
%
%  OPTIONAL, if a fourth image is wanted:
%      testbed_hardware.jpg — a photograph of the laptop, the USB gigabit
%      adapter, the single Ethernet cable and the Pi 4B. This is the only
%      image that shows the work is hardware-in-the-loop rather than pure
%      simulation. It would go in Section \ref{sec:test-env}.
% =====================================================================

%
%  Intended to follow % =====================================================================
%  Hardware-in-the-Loop Edge Processing
%  A SECTION of the FYP report. Include from the parent chapter with:
%      % =====================================================================
%  Hardware-in-the-Loop Edge Processing
%  A SECTION of the FYP report. Include from the parent chapter with:
%      % =====================================================================
%  Hardware-in-the-Loop Edge Processing
%  A SECTION of the FYP report. Include from the parent chapter with:
%      \input{hitl_edge_processing}
%
%  HEADING LEVELS USED
%      \section     - the HITL section itself
%      \subsection  - its parts
%      \paragraph   - short findings within Results
%
%  REQUIRED PACKAGES (add to the main preamble if not already present)
%      \usepackage{booktabs}
%      \usepackage{graphicx}
%      \usepackage{tikz}
%      \usetikzlibrary{arrows.meta, positioning, fit, backgrounds}
%
%  All labels are prefixed hitl / tab:hitl- / fig:hitl- so they cannot
%  collide with labels elsewhere in the shared report.
% =====================================================================

\section{Hardware-in-the-Loop Edge Processing}
\label{sec:hitl}

The vision workload was moved off the simulation host and onto physical
hardware, so that the cost of onboard detection could be measured on a processor
representative of one that would actually fly. A Raspberry~Pi~4B performs
detection on imagery produced by the simulation, and returns its results across a
wireless channel simulated in ns-3.

% ---------------------------------------------------------------------
\subsection{Extending the testbed with physical hardware}
\label{sec:hitl-testbed}

The existing stack ran entirely on one machine: Gazebo for the environment and
sensors, ArduPilot SITL for the autopilot, ns-3 for the wireless channel, and
Linux network namespaces isolating the ground station from each aircraft. The
ns-3 scenario defined four nodes --- a ground control station and three aircraft
--- of which the second and third were unused placeholders.

A Raspberry~Pi~4B was introduced as the physical edge node and assigned to the
second aircraft slot. Its TAP device, previously created only to satisfy the
scenario's argument list, was brought into service as the attachment point. The
autopilot, physics and channel simulation remained on the host; only the vision
computation moved onto hardware.

The Pi was configured with Ubuntu~22.04 and ROS\,2 Humble to match the host
exactly, giving both machines the same DDS implementation. The host has no
Ethernet port, so a USB~3.0 gigabit adapter was fitted. An initial low-cost
adapter was replaced after \texttt{iperf3} measured \(4.5\)\,Mbps with \(24\%\)
loss, despite \texttt{ethtool} reporting \(100\)\,Mbps full duplex; the
replacement sustained \(834\)\,Mbps with no retransmissions.

% ---------------------------------------------------------------------
\subsection{Constructing the two-link network}
\label{sec:hitl-network}

Two logically separate paths were required between host and Pi. One carries
camera imagery and must remain unimpaired, since on a real aircraft it is a short
cable inside a single airframe rather than a radio. The other carries detection
results and must pass through the simulated channel, as that is the link the
experiment exists to characterise. The same reasoning applies to the
flight-dynamics traffic between SITL and Gazebo, which also bypasses ns-3.

The Pi provides a single Ethernet interface, so both paths were carried over one
cable using IEEE~802.1Q VLAN tagging. On the host, two virtual interfaces were
created on the physical adapter: \texttt{eth-cam}, tagged VLAN~10 and addressed
\(10.0.0.1\); and \texttt{eth-rf}, tagged VLAN~42, left without an address and
made a member of a new bridge, \texttt{br-uav2}. That bridge also received
\texttt{tap-uav2}, the ns-3 interface for the second aircraft node. On the Pi,
matching interfaces \texttt{eth0.10} (\(10.0.0.2\)) and \texttt{eth0.42}
(\(10.42.0.12\)) were created.

Attaching the physical interface to the bridge directly was rejected during
design. A Linux bridge floods frames to all member ports, so camera traffic would
have been delivered into the ns-3 TAP device as well, injecting approximately
\(111\)\,Mbps of imagery into the simulated channel.

The resulting arrangement enforces separation through routing rather than
filtering. Gazebo is addressed only on \(10.0.0.1\) and can reach the Pi only at
\(10.0.0.2\); the ground station namespace is addressed only on \(10.42.0.10\)
and can reach it only at \(10.42.0.12\). Neither can use the other's path.

\begin{figure}[htbp]
\centering
\begin{tikzpicture}[
  font=\small,
  box/.style={draw, rounded corners=2pt, align=center, inner sep=5pt},
  proc/.style={box, fill=black!4},
  net/.style={box, fill=black!10},
  arr/.style={-{Latex[length=2mm]}, thick},
]
\node[proc] (gz)   {Gazebo\\camera + physics};
\node[proc, below=14mm of gz] (gcs) {Ground station\\(namespace)\\\texttt{10.42.0.10}};
\node[net,  right=16mm of gcs] (ns3) {ns-3\\channel};
\node[net,  above=8mm of ns3]  (br)  {\texttt{br-uav2}\\+ TAP};
\node[proc, right=46mm of gz] (pi) {Raspberry Pi 4B\\YOLOv8n detector};

\draw[arr] (gz) -- node[above, align=center, font=\scriptsize]
  {VLAN 10 \quad camera, 2.76\,MB/frame\\unimpaired} (pi);
\draw[arr] (pi.south) |- node[pos=0.75, above, font=\scriptsize]
  {VLAN 42 \quad detections, 118\,B} (br.east);
\draw[arr] (br) -- (ns3);
\draw[arr] (ns3) -- (gcs);

\begin{scope}[on background layer]
  \node[draw, dashed, rounded corners, fit=(gz)(gcs)(ns3)(br),
        inner sep=7pt, label={[font=\scriptsize\itshape]above left:Host PC}] {};
\end{scope}
\end{tikzpicture}
\caption{The two-link arrangement. Camera imagery reaches the edge node over an
unimpaired wired path; detection results traverse the simulated channel. Both
share one physical cable, separated by VLAN tagging.}
\label{fig:hitl-architecture}
\end{figure}

Addressing was made persistent using NetworkManager on the host and netplan on
the Pi, the two distributions using different network managers. Addresses
assigned imperatively were found to be removed within minutes, as each manager
reconciles kernel state against its own configuration and deletes entries it did
not create.

% ---------------------------------------------------------------------
\subsection{Configuring the middleware}
\label{sec:hitl-dds}

A Fast DDS profile was written and applied on both machines. It pins participants
to the intended interfaces through an interface whitelist, since both machines
are attached to WiFi as well as the wired link and DDS otherwise advertises on
every available interface. The profile also disables the shared-memory transport,
which is not isolated by network namespaces and would allow participants to
exchange data without traversing the simulated channel.

Kernel socket buffers were raised on both machines from the \(208\)\,kB default
to \(512\)\,MB. A single \(1280 \times 720\) frame is \(2.76\)\,MB and fragments
into approximately \(1{,}900\) UDP datagrams, of which the default buffer holds
about \(7\%\); reassembly never completed and subscribers received nothing, with
no error reported.

The camera subscription on the Pi was changed from the ROS\,2 default of
\texttt{RELIABLE} with queue depth ten to \texttt{BEST\_EFFORT} with depth one.
Under the default policy the publisher retains each sample until the subscriber
acknowledges it, and because the Pi acknowledges at approximately \(1\)\,Hz, the
camera output of the entire simulation was throttled to \(0.27\)\,Hz.
\texttt{RELIABLE} was retained for detection messages, which must not be lost.

% ---------------------------------------------------------------------
\subsection{Deploying the detector}
\label{sec:hitl-detector}

The detection node runs YOLOv8n restricted to the COCO \emph{person} class,
subscribing to the camera topic and publishing results as JSON. It was deployed
to the Pi and instrumented to record inference duration, measured around the
model invocation alone so that figures remain comparable between back-ends and
against standalone execution.

Two dependency conflicts were resolved during deployment. Installing
\texttt{ultralytics} brings in NumPy~2.x, against which \texttt{cv\_bridge} --- a
compiled extension built for the 1.x C ABI --- fails. Installing an older OpenCV
to compensate reintroduces NumPy~2 as a dependency, and OpenCV~5.0 separately
breaks \texttt{cv\_bridge} by renaming the type constants it uses to build its
conversion table. Versions were pinned to \texttt{numpy==1.26.4} and
\texttt{opencv-python==4.10.0.84}, the only combination satisfying both.

The model was additionally exported to NCNN for ARM-optimised inference. Because
NCNN fixes the input tensor shape at export time, the export geometry was
specified explicitly as \(544 \times 960\) to match the letterboxed frame; a
default export produces a square input in which approximately \(43\%\) of the
computation is spent on padding, at \(19.9\)\,GFLOPS against \(11.3\)\,GFLOPS for
the correct shape. Half-precision export was avoided, the Cortex-A72 lacking fp16
SIMD support.

% ---------------------------------------------------------------------
\subsection{Automation}
\label{sec:hitl-automation}

Bring-up was consolidated into a single script that starts the namespaces, ns-3,
Gazebo, SITL, the micro-ROS agent and the ground station receiver, creates the
VLANs and bridge, and launches the detector on the Pi over SSH. Pre-flight checks
were added for each failure encountered during development: presence of the host
address before Gazebo starts, since Fast DDS enumerates interfaces only at
participant creation and cannot use one assigned later; reachability of the Pi;
adequate socket buffer limits on both machines; and confirmation after startup
that Gazebo bound the intended interface.

% ---------------------------------------------------------------------
\subsection{Results}
\label{sec:hitl-results}

Detections were verified visually, by inspecting annotated frames, before any
quantitative result was accepted. An earlier configuration had reported up to
sixteen detections against five subjects; these were false positives on building
texture, eliminated by raising the confidence threshold and narrowing the camera
field of view.

\paragraph{Channel characterisation.}
Table~\ref{tab:hitl-links} compares the two paths, each measured over thirty ICMP
samples. Short samples are unreliable: the first packet includes address
resolution across the simulated channel, and a two-sample measurement of the same
configuration returned a mean of \(167\)\,ms against \(59\)\,ms for thirty.

\begin{table}[htbp]
\centering
\caption{Measured characteristics of the two links, thirty samples each.}
\label{tab:hitl-links}
\begin{tabular}{lrrrr}
\toprule
Path & Min (ms) & Mean (ms) & Max (ms) & Jitter (ms) \\
\midrule
Sensor, VLAN 10 (wired)   & 0.4  & 1.5  & 2.7   & 0.4  \\
Radio, VLAN 42 (via ns-3) & 11.7 & 77.4 & 242.6 & 54.8 \\
\bottomrule
\end{tabular}
\end{table}

Packet loss was zero on both paths, the aircraft remaining close to the ground
station throughout. The radio path is approximately \(53\times\) slower than the
wired path on average and varies approximately \(140\times\) more.

The distribution within the radio path is itself informative. Its minimum of
\(11.7\)\,ms against a mean of \(77.4\)\,ms is a spread of some \(6.6\times\) on
one link, where a channel applying a fixed delay would produce a narrow one.
The variation indicates that ns-3 is modelling contention, in which packets
queue and back off before transmission, rather than simply adding a constant
offset to each packet.

\paragraph{Sensor transport.}
With the camera at \(20\)\,Hz the wired path carried \(481\)\,Mbps at
\(42{,}544\) packets per second with zero dropped packets, occupying
approximately \(58\%\) of measured link capacity.

\paragraph{Inference latency.}
Table~\ref{tab:hitl-inference} reports latency for each combination of back-end
and input resolution. Changing back-end at constant resolution gives
approximately \(2.5\times\); halving the pixel count at constant back-end gives
approximately \(2.3\times\). Latency therefore scales close to linearly with
pixel count, a pixel ratio of \(0.47\) producing a time ratio of \(0.43\).

\begin{table}[htbp]
\centering
\caption{Inference latency on the Raspberry Pi 4B.}
\label{tab:hitl-inference}
\begin{tabular}{llrr}
\toprule
Back-end & Input & Latency (ms) & Rate (fps) \\
\midrule
PyTorch & $960 \times 544$ & 2540 & 0.39 \\
PyTorch & $640 \times 384$ & 1100 & 0.91 \\
NCNN    & $960 \times 544$ & 1000 & 1.00 \\
NCNN    & $640 \times 384$ & \multicolumn{2}{l}{not measured (est.\ 470)} \\
\bottomrule
\end{tabular}
\end{table}

\paragraph{Detection performance.}
Between two and four of the five subjects were detected per frame, a recall of
\(40\)--\(80\%\). The shortfall is attributable to scene geometry rather than
detector deficiency: the models are separated by approximately one metre, so from
the oblique camera view the nearer figures partially occlude those behind them,
and partially occluded subjects at approximately \(53\) pixels in height fall
below the confidence threshold. Because the number and position of subjects are
known exactly from the world definition, recall is a directly measurable quantity
rather than an estimate.

\paragraph{Transmitted volume.}
Each detection message is approximately \(118\)\,B, against \(2.76\)\,MB for the
corresponding frame, a reduction of approximately \(23{,}000\times\) in data
crossing the radio.

% ---------------------------------------------------------------------
\subsection{Contention between sensor transport and inference}
\label{sec:hitl-contention}

The two workloads on the edge node proved not to be independent. The same model,
on the same input, required \(623\)\,ms executed on a static image but
\(1{,}343\)\,ms within the live pipeline: a difference of \(720\)\,ms, or
approximately half the available inference throughput.

The edge node received frames at \(19.8\)\,Hz while completing inference at
approximately \(0.74\)\,Hz, discarding roughly \(96\%\) of what it received.
Discarded frames are not free: each requires reassembly of some \(2{,}150\) UDP
fragments, deserialisation of \(2.76\)\,MB and a further \(2.76\)\,MB
colour-space conversion, all competing with inference for the same four
processor cores.

Reducing the camera rate at source from \(20\)\,Hz to \(5\)\,Hz --- still
approximately seven times the rate the detector can consume --- reduced inference
latency to approximately \(1{,}000\)\,ms, recovering about a quarter of the
deficit, and reduced link occupancy from \(481\)\,Mbps to \(111\)\,Mbps. The
application-level frame-rate control provided by the relay node cannot achieve
this, as it discards frames after reception, by which point the transport and
deserialisation cost has already been incurred.

This interaction is not observable in pure simulation, where sensor data is
delivered through shared memory at negligible cost, nor in an isolated benchmark
of the detector, since the interference originates outside it. It appears only
when a physical processor performs both reception and inference concurrently.

% ---------------------------------------------------------------------
\subsection{Limitations}
\label{sec:hitl-limitations}

Clock synchronisation between the two machines has not yet been established.
End-to-end latency is computed from timestamps taken on different hosts, and
without synchronisation those values are not meaningful while remaining plausible
in magnitude; synchronisation is required before end-to-end latency is reported.

The edge configuration is characterised; the corresponding ground configuration,
in which compressed imagery crosses the simulated channel and inference is
performed at the ground station, remains to be measured.

Measurements were taken with the aircraft close to the ground station, giving
zero packet loss throughout. Behaviour as range increases and buildings obstruct
the path has not been characterised. The testbed supports a single physical edge
node; multi-node operation has not been evaluated.

%
%  HEADING LEVELS USED
%      \section     - the HITL section itself
%      \subsection  - its parts
%      \paragraph   - short findings within Results
%
%  REQUIRED PACKAGES (add to the main preamble if not already present)
%      \usepackage{booktabs}
%      \usepackage{graphicx}
%      \usepackage{tikz}
%      \usetikzlibrary{arrows.meta, positioning, fit, backgrounds}
%
%  All labels are prefixed hitl / tab:hitl- / fig:hitl- so they cannot
%  collide with labels elsewhere in the shared report.
% =====================================================================

\section{Hardware-in-the-Loop Edge Processing}
\label{sec:hitl}

The vision workload was moved off the simulation host and onto physical
hardware, so that the cost of onboard detection could be measured on a processor
representative of one that would actually fly. A Raspberry~Pi~4B performs
detection on imagery produced by the simulation, and returns its results across a
wireless channel simulated in ns-3.

% ---------------------------------------------------------------------
\subsection{Extending the testbed with physical hardware}
\label{sec:hitl-testbed}

The existing stack ran entirely on one machine: Gazebo for the environment and
sensors, ArduPilot SITL for the autopilot, ns-3 for the wireless channel, and
Linux network namespaces isolating the ground station from each aircraft. The
ns-3 scenario defined four nodes --- a ground control station and three aircraft
--- of which the second and third were unused placeholders.

A Raspberry~Pi~4B was introduced as the physical edge node and assigned to the
second aircraft slot. Its TAP device, previously created only to satisfy the
scenario's argument list, was brought into service as the attachment point. The
autopilot, physics and channel simulation remained on the host; only the vision
computation moved onto hardware.

The Pi was configured with Ubuntu~22.04 and ROS\,2 Humble to match the host
exactly, giving both machines the same DDS implementation. The host has no
Ethernet port, so a USB~3.0 gigabit adapter was fitted. An initial low-cost
adapter was replaced after \texttt{iperf3} measured \(4.5\)\,Mbps with \(24\%\)
loss, despite \texttt{ethtool} reporting \(100\)\,Mbps full duplex; the
replacement sustained \(834\)\,Mbps with no retransmissions.

% ---------------------------------------------------------------------
\subsection{Constructing the two-link network}
\label{sec:hitl-network}

Two logically separate paths were required between host and Pi. One carries
camera imagery and must remain unimpaired, since on a real aircraft it is a short
cable inside a single airframe rather than a radio. The other carries detection
results and must pass through the simulated channel, as that is the link the
experiment exists to characterise. The same reasoning applies to the
flight-dynamics traffic between SITL and Gazebo, which also bypasses ns-3.

The Pi provides a single Ethernet interface, so both paths were carried over one
cable using IEEE~802.1Q VLAN tagging. On the host, two virtual interfaces were
created on the physical adapter: \texttt{eth-cam}, tagged VLAN~10 and addressed
\(10.0.0.1\); and \texttt{eth-rf}, tagged VLAN~42, left without an address and
made a member of a new bridge, \texttt{br-uav2}. That bridge also received
\texttt{tap-uav2}, the ns-3 interface for the second aircraft node. On the Pi,
matching interfaces \texttt{eth0.10} (\(10.0.0.2\)) and \texttt{eth0.42}
(\(10.42.0.12\)) were created.

Attaching the physical interface to the bridge directly was rejected during
design. A Linux bridge floods frames to all member ports, so camera traffic would
have been delivered into the ns-3 TAP device as well, injecting approximately
\(111\)\,Mbps of imagery into the simulated channel.

The resulting arrangement enforces separation through routing rather than
filtering. Gazebo is addressed only on \(10.0.0.1\) and can reach the Pi only at
\(10.0.0.2\); the ground station namespace is addressed only on \(10.42.0.10\)
and can reach it only at \(10.42.0.12\). Neither can use the other's path.

\begin{figure}[htbp]
\centering
\begin{tikzpicture}[
  font=\small,
  box/.style={draw, rounded corners=2pt, align=center, inner sep=5pt},
  proc/.style={box, fill=black!4},
  net/.style={box, fill=black!10},
  arr/.style={-{Latex[length=2mm]}, thick},
]
\node[proc] (gz)   {Gazebo\\camera + physics};
\node[proc, below=14mm of gz] (gcs) {Ground station\\(namespace)\\\texttt{10.42.0.10}};
\node[net,  right=16mm of gcs] (ns3) {ns-3\\channel};
\node[net,  above=8mm of ns3]  (br)  {\texttt{br-uav2}\\+ TAP};
\node[proc, right=46mm of gz] (pi) {Raspberry Pi 4B\\YOLOv8n detector};

\draw[arr] (gz) -- node[above, align=center, font=\scriptsize]
  {VLAN 10 \quad camera, 2.76\,MB/frame\\unimpaired} (pi);
\draw[arr] (pi.south) |- node[pos=0.75, above, font=\scriptsize]
  {VLAN 42 \quad detections, 118\,B} (br.east);
\draw[arr] (br) -- (ns3);
\draw[arr] (ns3) -- (gcs);

\begin{scope}[on background layer]
  \node[draw, dashed, rounded corners, fit=(gz)(gcs)(ns3)(br),
        inner sep=7pt, label={[font=\scriptsize\itshape]above left:Host PC}] {};
\end{scope}
\end{tikzpicture}
\caption{The two-link arrangement. Camera imagery reaches the edge node over an
unimpaired wired path; detection results traverse the simulated channel. Both
share one physical cable, separated by VLAN tagging.}
\label{fig:hitl-architecture}
\end{figure}

Addressing was made persistent using NetworkManager on the host and netplan on
the Pi, the two distributions using different network managers. Addresses
assigned imperatively were found to be removed within minutes, as each manager
reconciles kernel state against its own configuration and deletes entries it did
not create.

% ---------------------------------------------------------------------
\subsection{Configuring the middleware}
\label{sec:hitl-dds}

A Fast DDS profile was written and applied on both machines. It pins participants
to the intended interfaces through an interface whitelist, since both machines
are attached to WiFi as well as the wired link and DDS otherwise advertises on
every available interface. The profile also disables the shared-memory transport,
which is not isolated by network namespaces and would allow participants to
exchange data without traversing the simulated channel.

Kernel socket buffers were raised on both machines from the \(208\)\,kB default
to \(512\)\,MB. A single \(1280 \times 720\) frame is \(2.76\)\,MB and fragments
into approximately \(1{,}900\) UDP datagrams, of which the default buffer holds
about \(7\%\); reassembly never completed and subscribers received nothing, with
no error reported.

The camera subscription on the Pi was changed from the ROS\,2 default of
\texttt{RELIABLE} with queue depth ten to \texttt{BEST\_EFFORT} with depth one.
Under the default policy the publisher retains each sample until the subscriber
acknowledges it, and because the Pi acknowledges at approximately \(1\)\,Hz, the
camera output of the entire simulation was throttled to \(0.27\)\,Hz.
\texttt{RELIABLE} was retained for detection messages, which must not be lost.

% ---------------------------------------------------------------------
\subsection{Deploying the detector}
\label{sec:hitl-detector}

The detection node runs YOLOv8n restricted to the COCO \emph{person} class,
subscribing to the camera topic and publishing results as JSON. It was deployed
to the Pi and instrumented to record inference duration, measured around the
model invocation alone so that figures remain comparable between back-ends and
against standalone execution.

Two dependency conflicts were resolved during deployment. Installing
\texttt{ultralytics} brings in NumPy~2.x, against which \texttt{cv\_bridge} --- a
compiled extension built for the 1.x C ABI --- fails. Installing an older OpenCV
to compensate reintroduces NumPy~2 as a dependency, and OpenCV~5.0 separately
breaks \texttt{cv\_bridge} by renaming the type constants it uses to build its
conversion table. Versions were pinned to \texttt{numpy==1.26.4} and
\texttt{opencv-python==4.10.0.84}, the only combination satisfying both.

The model was additionally exported to NCNN for ARM-optimised inference. Because
NCNN fixes the input tensor shape at export time, the export geometry was
specified explicitly as \(544 \times 960\) to match the letterboxed frame; a
default export produces a square input in which approximately \(43\%\) of the
computation is spent on padding, at \(19.9\)\,GFLOPS against \(11.3\)\,GFLOPS for
the correct shape. Half-precision export was avoided, the Cortex-A72 lacking fp16
SIMD support.

% ---------------------------------------------------------------------
\subsection{Automation}
\label{sec:hitl-automation}

Bring-up was consolidated into a single script that starts the namespaces, ns-3,
Gazebo, SITL, the micro-ROS agent and the ground station receiver, creates the
VLANs and bridge, and launches the detector on the Pi over SSH. Pre-flight checks
were added for each failure encountered during development: presence of the host
address before Gazebo starts, since Fast DDS enumerates interfaces only at
participant creation and cannot use one assigned later; reachability of the Pi;
adequate socket buffer limits on both machines; and confirmation after startup
that Gazebo bound the intended interface.

% ---------------------------------------------------------------------
\subsection{Results}
\label{sec:hitl-results}

Detections were verified visually, by inspecting annotated frames, before any
quantitative result was accepted. An earlier configuration had reported up to
sixteen detections against five subjects; these were false positives on building
texture, eliminated by raising the confidence threshold and narrowing the camera
field of view.

\paragraph{Channel characterisation.}
Table~\ref{tab:hitl-links} compares the two paths, each measured over thirty ICMP
samples. Short samples are unreliable: the first packet includes address
resolution across the simulated channel, and a two-sample measurement of the same
configuration returned a mean of \(167\)\,ms against \(59\)\,ms for thirty.

\begin{table}[htbp]
\centering
\caption{Measured characteristics of the two links, thirty samples each.}
\label{tab:hitl-links}
\begin{tabular}{lrrrr}
\toprule
Path & Min (ms) & Mean (ms) & Max (ms) & Jitter (ms) \\
\midrule
Sensor, VLAN 10 (wired)   & 0.4  & 1.5  & 2.7   & 0.4  \\
Radio, VLAN 42 (via ns-3) & 11.7 & 77.4 & 242.6 & 54.8 \\
\bottomrule
\end{tabular}
\end{table}

Packet loss was zero on both paths, the aircraft remaining close to the ground
station throughout. The radio path is approximately \(53\times\) slower than the
wired path on average and varies approximately \(140\times\) more.

The distribution within the radio path is itself informative. Its minimum of
\(11.7\)\,ms against a mean of \(77.4\)\,ms is a spread of some \(6.6\times\) on
one link, where a channel applying a fixed delay would produce a narrow one.
The variation indicates that ns-3 is modelling contention, in which packets
queue and back off before transmission, rather than simply adding a constant
offset to each packet.

\paragraph{Sensor transport.}
With the camera at \(20\)\,Hz the wired path carried \(481\)\,Mbps at
\(42{,}544\) packets per second with zero dropped packets, occupying
approximately \(58\%\) of measured link capacity.

\paragraph{Inference latency.}
Table~\ref{tab:hitl-inference} reports latency for each combination of back-end
and input resolution. Changing back-end at constant resolution gives
approximately \(2.5\times\); halving the pixel count at constant back-end gives
approximately \(2.3\times\). Latency therefore scales close to linearly with
pixel count, a pixel ratio of \(0.47\) producing a time ratio of \(0.43\).

\begin{table}[htbp]
\centering
\caption{Inference latency on the Raspberry Pi 4B.}
\label{tab:hitl-inference}
\begin{tabular}{llrr}
\toprule
Back-end & Input & Latency (ms) & Rate (fps) \\
\midrule
PyTorch & $960 \times 544$ & 2540 & 0.39 \\
PyTorch & $640 \times 384$ & 1100 & 0.91 \\
NCNN    & $960 \times 544$ & 1000 & 1.00 \\
NCNN    & $640 \times 384$ & \multicolumn{2}{l}{not measured (est.\ 470)} \\
\bottomrule
\end{tabular}
\end{table}

\paragraph{Detection performance.}
Between two and four of the five subjects were detected per frame, a recall of
\(40\)--\(80\%\). The shortfall is attributable to scene geometry rather than
detector deficiency: the models are separated by approximately one metre, so from
the oblique camera view the nearer figures partially occlude those behind them,
and partially occluded subjects at approximately \(53\) pixels in height fall
below the confidence threshold. Because the number and position of subjects are
known exactly from the world definition, recall is a directly measurable quantity
rather than an estimate.

\paragraph{Transmitted volume.}
Each detection message is approximately \(118\)\,B, against \(2.76\)\,MB for the
corresponding frame, a reduction of approximately \(23{,}000\times\) in data
crossing the radio.

% ---------------------------------------------------------------------
\subsection{Contention between sensor transport and inference}
\label{sec:hitl-contention}

The two workloads on the edge node proved not to be independent. The same model,
on the same input, required \(623\)\,ms executed on a static image but
\(1{,}343\)\,ms within the live pipeline: a difference of \(720\)\,ms, or
approximately half the available inference throughput.

The edge node received frames at \(19.8\)\,Hz while completing inference at
approximately \(0.74\)\,Hz, discarding roughly \(96\%\) of what it received.
Discarded frames are not free: each requires reassembly of some \(2{,}150\) UDP
fragments, deserialisation of \(2.76\)\,MB and a further \(2.76\)\,MB
colour-space conversion, all competing with inference for the same four
processor cores.

Reducing the camera rate at source from \(20\)\,Hz to \(5\)\,Hz --- still
approximately seven times the rate the detector can consume --- reduced inference
latency to approximately \(1{,}000\)\,ms, recovering about a quarter of the
deficit, and reduced link occupancy from \(481\)\,Mbps to \(111\)\,Mbps. The
application-level frame-rate control provided by the relay node cannot achieve
this, as it discards frames after reception, by which point the transport and
deserialisation cost has already been incurred.

This interaction is not observable in pure simulation, where sensor data is
delivered through shared memory at negligible cost, nor in an isolated benchmark
of the detector, since the interference originates outside it. It appears only
when a physical processor performs both reception and inference concurrently.

% ---------------------------------------------------------------------
\subsection{Limitations}
\label{sec:hitl-limitations}

Clock synchronisation between the two machines has not yet been established.
End-to-end latency is computed from timestamps taken on different hosts, and
without synchronisation those values are not meaningful while remaining plausible
in magnitude; synchronisation is required before end-to-end latency is reported.

The edge configuration is characterised; the corresponding ground configuration,
in which compressed imagery crosses the simulated channel and inference is
performed at the ground station, remains to be measured.

Measurements were taken with the aircraft close to the ground station, giving
zero packet loss throughout. Behaviour as range increases and buildings obstruct
the path has not been characterised. The testbed supports a single physical edge
node; multi-node operation has not been evaluated.

%
%  HEADING LEVELS USED
%      \section     - the HITL section itself
%      \subsection  - its parts
%      \paragraph   - short findings within Results
%
%  REQUIRED PACKAGES (add to the main preamble if not already present)
%      \usepackage{booktabs}
%      \usepackage{graphicx}
%      \usepackage{tikz}
%      \usetikzlibrary{arrows.meta, positioning, fit, backgrounds}
%
%  All labels are prefixed hitl / tab:hitl- / fig:hitl- so they cannot
%  collide with labels elsewhere in the shared report.
% =====================================================================

\section{Hardware-in-the-Loop Edge Processing}
\label{sec:hitl}

The vision workload was moved off the simulation host and onto physical
hardware, so that the cost of onboard detection could be measured on a processor
representative of one that would actually fly. A Raspberry~Pi~4B performs
detection on imagery produced by the simulation, and returns its results across a
wireless channel simulated in ns-3.

% ---------------------------------------------------------------------
\subsection{Extending the testbed with physical hardware}
\label{sec:hitl-testbed}

The existing stack ran entirely on one machine: Gazebo for the environment and
sensors, ArduPilot SITL for the autopilot, ns-3 for the wireless channel, and
Linux network namespaces isolating the ground station from each aircraft. The
ns-3 scenario defined four nodes --- a ground control station and three aircraft
--- of which the second and third were unused placeholders.

A Raspberry~Pi~4B was introduced as the physical edge node and assigned to the
second aircraft slot. Its TAP device, previously created only to satisfy the
scenario's argument list, was brought into service as the attachment point. The
autopilot, physics and channel simulation remained on the host; only the vision
computation moved onto hardware.

The Pi was configured with Ubuntu~22.04 and ROS\,2 Humble to match the host
exactly, giving both machines the same DDS implementation. The host has no
Ethernet port, so a USB~3.0 gigabit adapter was fitted. An initial low-cost
adapter was replaced after \texttt{iperf3} measured \(4.5\)\,Mbps with \(24\%\)
loss, despite \texttt{ethtool} reporting \(100\)\,Mbps full duplex; the
replacement sustained \(834\)\,Mbps with no retransmissions.

% ---------------------------------------------------------------------
\subsection{Constructing the two-link network}
\label{sec:hitl-network}

Two logically separate paths were required between host and Pi. One carries
camera imagery and must remain unimpaired, since on a real aircraft it is a short
cable inside a single airframe rather than a radio. The other carries detection
results and must pass through the simulated channel, as that is the link the
experiment exists to characterise. The same reasoning applies to the
flight-dynamics traffic between SITL and Gazebo, which also bypasses ns-3.

The Pi provides a single Ethernet interface, so both paths were carried over one
cable using IEEE~802.1Q VLAN tagging. On the host, two virtual interfaces were
created on the physical adapter: \texttt{eth-cam}, tagged VLAN~10 and addressed
\(10.0.0.1\); and \texttt{eth-rf}, tagged VLAN~42, left without an address and
made a member of a new bridge, \texttt{br-uav2}. That bridge also received
\texttt{tap-uav2}, the ns-3 interface for the second aircraft node. On the Pi,
matching interfaces \texttt{eth0.10} (\(10.0.0.2\)) and \texttt{eth0.42}
(\(10.42.0.12\)) were created.

Attaching the physical interface to the bridge directly was rejected during
design. A Linux bridge floods frames to all member ports, so camera traffic would
have been delivered into the ns-3 TAP device as well, injecting approximately
\(111\)\,Mbps of imagery into the simulated channel.

The resulting arrangement enforces separation through routing rather than
filtering. Gazebo is addressed only on \(10.0.0.1\) and can reach the Pi only at
\(10.0.0.2\); the ground station namespace is addressed only on \(10.42.0.10\)
and can reach it only at \(10.42.0.12\). Neither can use the other's path.

\begin{figure}[htbp]
\centering
\begin{tikzpicture}[
  font=\small,
  box/.style={draw, rounded corners=2pt, align=center, inner sep=5pt},
  proc/.style={box, fill=black!4},
  net/.style={box, fill=black!10},
  arr/.style={-{Latex[length=2mm]}, thick},
]
\node[proc] (gz)   {Gazebo\\camera + physics};
\node[proc, below=14mm of gz] (gcs) {Ground station\\(namespace)\\\texttt{10.42.0.10}};
\node[net,  right=16mm of gcs] (ns3) {ns-3\\channel};
\node[net,  above=8mm of ns3]  (br)  {\texttt{br-uav2}\\+ TAP};
\node[proc, right=46mm of gz] (pi) {Raspberry Pi 4B\\YOLOv8n detector};

\draw[arr] (gz) -- node[above, align=center, font=\scriptsize]
  {VLAN 10 \quad camera, 2.76\,MB/frame\\unimpaired} (pi);
\draw[arr] (pi.south) |- node[pos=0.75, above, font=\scriptsize]
  {VLAN 42 \quad detections, 118\,B} (br.east);
\draw[arr] (br) -- (ns3);
\draw[arr] (ns3) -- (gcs);

\begin{scope}[on background layer]
  \node[draw, dashed, rounded corners, fit=(gz)(gcs)(ns3)(br),
        inner sep=7pt, label={[font=\scriptsize\itshape]above left:Host PC}] {};
\end{scope}
\end{tikzpicture}
\caption{The two-link arrangement. Camera imagery reaches the edge node over an
unimpaired wired path; detection results traverse the simulated channel. Both
share one physical cable, separated by VLAN tagging.}
\label{fig:hitl-architecture}
\end{figure}

Addressing was made persistent using NetworkManager on the host and netplan on
the Pi, the two distributions using different network managers. Addresses
assigned imperatively were found to be removed within minutes, as each manager
reconciles kernel state against its own configuration and deletes entries it did
not create.

% ---------------------------------------------------------------------
\subsection{Configuring the middleware}
\label{sec:hitl-dds}

A Fast DDS profile was written and applied on both machines. It pins participants
to the intended interfaces through an interface whitelist, since both machines
are attached to WiFi as well as the wired link and DDS otherwise advertises on
every available interface. The profile also disables the shared-memory transport,
which is not isolated by network namespaces and would allow participants to
exchange data without traversing the simulated channel.

Kernel socket buffers were raised on both machines from the \(208\)\,kB default
to \(512\)\,MB. A single \(1280 \times 720\) frame is \(2.76\)\,MB and fragments
into approximately \(1{,}900\) UDP datagrams, of which the default buffer holds
about \(7\%\); reassembly never completed and subscribers received nothing, with
no error reported.

The camera subscription on the Pi was changed from the ROS\,2 default of
\texttt{RELIABLE} with queue depth ten to \texttt{BEST\_EFFORT} with depth one.
Under the default policy the publisher retains each sample until the subscriber
acknowledges it, and because the Pi acknowledges at approximately \(1\)\,Hz, the
camera output of the entire simulation was throttled to \(0.27\)\,Hz.
\texttt{RELIABLE} was retained for detection messages, which must not be lost.

% ---------------------------------------------------------------------
\subsection{Deploying the detector}
\label{sec:hitl-detector}

The detection node runs YOLOv8n restricted to the COCO \emph{person} class,
subscribing to the camera topic and publishing results as JSON. It was deployed
to the Pi and instrumented to record inference duration, measured around the
model invocation alone so that figures remain comparable between back-ends and
against standalone execution.

Two dependency conflicts were resolved during deployment. Installing
\texttt{ultralytics} brings in NumPy~2.x, against which \texttt{cv\_bridge} --- a
compiled extension built for the 1.x C ABI --- fails. Installing an older OpenCV
to compensate reintroduces NumPy~2 as a dependency, and OpenCV~5.0 separately
breaks \texttt{cv\_bridge} by renaming the type constants it uses to build its
conversion table. Versions were pinned to \texttt{numpy==1.26.4} and
\texttt{opencv-python==4.10.0.84}, the only combination satisfying both.

The model was additionally exported to NCNN for ARM-optimised inference. Because
NCNN fixes the input tensor shape at export time, the export geometry was
specified explicitly as \(544 \times 960\) to match the letterboxed frame; a
default export produces a square input in which approximately \(43\%\) of the
computation is spent on padding, at \(19.9\)\,GFLOPS against \(11.3\)\,GFLOPS for
the correct shape. Half-precision export was avoided, the Cortex-A72 lacking fp16
SIMD support.

% ---------------------------------------------------------------------
\subsection{Automation}
\label{sec:hitl-automation}

Bring-up was consolidated into a single script that starts the namespaces, ns-3,
Gazebo, SITL, the micro-ROS agent and the ground station receiver, creates the
VLANs and bridge, and launches the detector on the Pi over SSH. Pre-flight checks
were added for each failure encountered during development: presence of the host
address before Gazebo starts, since Fast DDS enumerates interfaces only at
participant creation and cannot use one assigned later; reachability of the Pi;
adequate socket buffer limits on both machines; and confirmation after startup
that Gazebo bound the intended interface.

% ---------------------------------------------------------------------
\subsection{Results}
\label{sec:hitl-results}

Detections were verified visually, by inspecting annotated frames, before any
quantitative result was accepted. An earlier configuration had reported up to
sixteen detections against five subjects; these were false positives on building
texture, eliminated by raising the confidence threshold and narrowing the camera
field of view.

\paragraph{Channel characterisation.}
Table~\ref{tab:hitl-links} compares the two paths, each measured over thirty ICMP
samples. Short samples are unreliable: the first packet includes address
resolution across the simulated channel, and a two-sample measurement of the same
configuration returned a mean of \(167\)\,ms against \(59\)\,ms for thirty.

\begin{table}[htbp]
\centering
\caption{Measured characteristics of the two links, thirty samples each.}
\label{tab:hitl-links}
\begin{tabular}{lrrrr}
\toprule
Path & Min (ms) & Mean (ms) & Max (ms) & Jitter (ms) \\
\midrule
Sensor, VLAN 10 (wired)   & 0.4  & 1.5  & 2.7   & 0.4  \\
Radio, VLAN 42 (via ns-3) & 11.7 & 77.4 & 242.6 & 54.8 \\
\bottomrule
\end{tabular}
\end{table}

Packet loss was zero on both paths, the aircraft remaining close to the ground
station throughout. The radio path is approximately \(53\times\) slower than the
wired path on average and varies approximately \(140\times\) more.

The distribution within the radio path is itself informative. Its minimum of
\(11.7\)\,ms against a mean of \(77.4\)\,ms is a spread of some \(6.6\times\) on
one link, where a channel applying a fixed delay would produce a narrow one.
The variation indicates that ns-3 is modelling contention, in which packets
queue and back off before transmission, rather than simply adding a constant
offset to each packet.

\paragraph{Sensor transport.}
With the camera at \(20\)\,Hz the wired path carried \(481\)\,Mbps at
\(42{,}544\) packets per second with zero dropped packets, occupying
approximately \(58\%\) of measured link capacity.

\paragraph{Inference latency.}
Table~\ref{tab:hitl-inference} reports latency for each combination of back-end
and input resolution. Changing back-end at constant resolution gives
approximately \(2.5\times\); halving the pixel count at constant back-end gives
approximately \(2.3\times\). Latency therefore scales close to linearly with
pixel count, a pixel ratio of \(0.47\) producing a time ratio of \(0.43\).

\begin{table}[htbp]
\centering
\caption{Inference latency on the Raspberry Pi 4B.}
\label{tab:hitl-inference}
\begin{tabular}{llrr}
\toprule
Back-end & Input & Latency (ms) & Rate (fps) \\
\midrule
PyTorch & $960 \times 544$ & 2540 & 0.39 \\
PyTorch & $640 \times 384$ & 1100 & 0.91 \\
NCNN    & $960 \times 544$ & 1000 & 1.00 \\
NCNN    & $640 \times 384$ & \multicolumn{2}{l}{not measured (est.\ 470)} \\
\bottomrule
\end{tabular}
\end{table}

\paragraph{Detection performance.}
Between two and four of the five subjects were detected per frame, a recall of
\(40\)--\(80\%\). The shortfall is attributable to scene geometry rather than
detector deficiency: the models are separated by approximately one metre, so from
the oblique camera view the nearer figures partially occlude those behind them,
and partially occluded subjects at approximately \(53\) pixels in height fall
below the confidence threshold. Because the number and position of subjects are
known exactly from the world definition, recall is a directly measurable quantity
rather than an estimate.

\paragraph{Transmitted volume.}
Each detection message is approximately \(118\)\,B, against \(2.76\)\,MB for the
corresponding frame, a reduction of approximately \(23{,}000\times\) in data
crossing the radio.

% ---------------------------------------------------------------------
\subsection{Contention between sensor transport and inference}
\label{sec:hitl-contention}

The two workloads on the edge node proved not to be independent. The same model,
on the same input, required \(623\)\,ms executed on a static image but
\(1{,}343\)\,ms within the live pipeline: a difference of \(720\)\,ms, or
approximately half the available inference throughput.

The edge node received frames at \(19.8\)\,Hz while completing inference at
approximately \(0.74\)\,Hz, discarding roughly \(96\%\) of what it received.
Discarded frames are not free: each requires reassembly of some \(2{,}150\) UDP
fragments, deserialisation of \(2.76\)\,MB and a further \(2.76\)\,MB
colour-space conversion, all competing with inference for the same four
processor cores.

Reducing the camera rate at source from \(20\)\,Hz to \(5\)\,Hz --- still
approximately seven times the rate the detector can consume --- reduced inference
latency to approximately \(1{,}000\)\,ms, recovering about a quarter of the
deficit, and reduced link occupancy from \(481\)\,Mbps to \(111\)\,Mbps. The
application-level frame-rate control provided by the relay node cannot achieve
this, as it discards frames after reception, by which point the transport and
deserialisation cost has already been incurred.

This interaction is not observable in pure simulation, where sensor data is
delivered through shared memory at negligible cost, nor in an isolated benchmark
of the detector, since the interference originates outside it. It appears only
when a physical processor performs both reception and inference concurrently.

% ---------------------------------------------------------------------
\subsection{Limitations}
\label{sec:hitl-limitations}

Clock synchronisation between the two machines has not yet been established.
End-to-end latency is computed from timestamps taken on different hosts, and
without synchronisation those values are not meaningful while remaining plausible
in magnitude; synchronisation is required before end-to-end latency is reported.

The edge configuration is characterised; the corresponding ground configuration,
in which compressed imagery crosses the simulated channel and inference is
performed at the ground station, remains to be measured.

Measurements were taken with the aircraft close to the ground station, giving
zero packet loss throughout. Behaviour as range increases and buildings obstruct
the path has not been characterised. The testbed supports a single physical edge
node; multi-node operation has not been evaluated.
; cross-references
%  \ref{sec:hitl}. A longer version with the full 26-case test matrix is
%  kept in testing_validation_long.tex if an appendix is wanted.
%
%  REQUIRED PACKAGES
%      \usepackage{booktabs}
%      \usepackage{array}
%      \usepackage{graphicx}
%      \usepackage{subcaption}
%      \graphicspath{{figures/}}
%
%  IMAGES EXPECTED in docs/report/figures/  (see checklist at end)
%      detection_a.png  detection_b.png  detection_c.png
% =====================================================================

\section{Testing, Experimentation and Validation}
\label{sec:test}

Testing of the hardware-in-the-loop testbed described in
Section~\ref{sec:hitl} was organised around a difficulty specific to this
kind of system: a fault at a low layer does not raise an error at the layer
above it, but produces plausible and wrong behaviour there. A subscriber
starved of buffer space receives nothing and reports nothing; two nodes with
mismatched topic names run indefinitely without complaint. Testing was
therefore performed bottom-up --- physical link, path separation, middleware
transport, detector, end-to-end --- with each layer required to pass before
the next was attempted.

Three results are reported here. The first establishes that the two-link
architecture does what it claims. The second establishes that the detector
finds real people rather than plausible-looking texture. The third concerns
the measurement chain itself, and is included because it determines whether
the other two can be believed.

% ---------------------------------------------------------------------
\subsection{Test environment}
\label{sec:test-env}

Both machines run Ubuntu~22.04 with ROS\,2 Humble and Fast~DDS~2.6, so the
middleware is identical at each end of the link and cannot itself be a
source of difference. The simulation host runs Gazebo Classic in lockstep
with ArduPilot SITL, and ns-3 with TapBridge for the wireless channel. The
edge node is a Raspberry~Pi~4B running YOLOv8n restricted to the COCO
\emph{person} class. The camera produces $1280 \times 720$ frames of
$2.76$\,MB at $5$\,Hz. The test scene contains exactly five human models at
coordinates fixed in the world definition, which makes recall a counted
quantity rather than an estimated one.

The Pi's single Ethernet port is divided by 802.1Q tagging into
\texttt{eth0.10} ($10.0.0.2$, camera, unimpaired) and \texttt{eth0.42}
($10.42.0.12$, detections, through ns-3), created by
\texttt{scripts/netns/pi\_hitl\_link.sh}. The matching host interfaces
\texttt{eth-cam} and \texttt{eth-rf} are created by step~1c of
\texttt{scripts/netns/launch\_single\_uav\_netns.sh}, the latter bridged
into \texttt{tap-uav2}.

% ---------------------------------------------------------------------
\subsection{Result 1: the two links behave differently}
\label{sec:test-links}

The central claim of the architecture is that detection results cross a
simulated wireless channel while camera imagery does not. Two paths that
both work are not evidence of this; they must also be shown to differ in the
intended way, and imagery must be shown to be absent from the radio path.

Table~\ref{tab:test-links} reports sixty ICMP samples on the sensor path and
thirty on the radio path. Large samples are necessary: the first packet on the
radio path includes address resolution across the simulated channel, and a
two-sample measurement of the same configuration returned a mean of
$167$\,ms against $74.5$\,ms for thirty. The short sample is biased, not
merely noisy.

\begin{table}[htbp]
\centering
\caption{Measured characteristics of the two paths. Sensor path over sixty
samples, radio path over thirty; zero packet loss on both.}
\label{tab:test-links}
\begin{tabular}{@{}lrrrr@{}}
\toprule
Path & Min (ms) & Mean (ms) & Max (ms) & Jitter (ms) \\
\midrule
Sensor, VLAN 10 (wired)   & 1.0  & 1.4  & 1.8   & 0.12 \\
Radio, VLAN 42 (via ns-3) & 10.1 & 74.5 & 158.0 & 38.4 \\
\midrule
Ratio                     & --   & $53\times$ & -- & $320\times$ \\
\bottomrule
\end{tabular}
\end{table}

The radio path is approximately \(53\times\) slower on average, and its
variation is larger by more than two orders of magnitude.

Its internal spread is itself informative: a minimum of \(10.1\)\,ms against
a mean of \(74.5\)\,ms is a range of some \(7\times\) on one link, where a
channel applying a fixed delay would produce a narrow distribution. The
variation indicates that ns-3 is modelling contention, in which packets queue
and back off, rather than adding a constant offset.

The radio measurement was repeated. An independent set of thirty samples taken
several hours earlier, with the host reconfigured in between, gave a mean of
\(77.4\)\,ms against the \(74.5\)\,ms reported here, agreeing within
\(4\%\).

The sensor figures were taken \emph{while} the camera was streaming at
\(110\)\,Mbps. An idle-cable measurement over the same number of samples
gives the same mean, so the two logical links sharing one physical cable do
not delay one another. Queue management was verified rather than assumed:
both machines run \texttt{fq\_codel} with a \(5\)\,ms target, and byte
queue limits on the edge node had reduced its driver ring to approximately one
packet.

Absence of imagery on the radio path was established negatively, by
observation rather than by inspecting a rule set, since separation here is
enforced by addressing and bridge membership rather than by filtering. No
$10.0.0.0/24$ traffic was observed on \texttt{tap-uav2}, the point at which
traffic enters ns-3, and the byte counters on the two interfaces differ by
approximately four orders of magnitude --- consistent with $2.76$\,MB per
frame on one path against $118$\,B per detection message on the other, a
reduction of approximately \(23{,}000\times\).

A structural argument supports the measurements. The ground station receiver
executes inside the \texttt{gcsns} network namespace, whose only interface
is one end of a veth pair attached to \texttt{tap-gcs}. That namespace has
no interface on $10.0.0.0/24$ and no route to it. A message arriving at the
receiver cannot have taken any other path, because no other path exists from
inside the namespace. Separation is therefore a property of the topology
rather than of a setting that could be silently overridden.

% ---------------------------------------------------------------------
\subsection{Result 2: detections correspond to real subjects}
\label{sec:test-detection}

Visual confirmation was made a precondition of accepting any quantitative
result. A detector reporting a plausible count while placing its boxes on
building texture would satisfy every numerical check in this section, and an
earlier configuration did exactly that: sixteen confident detections against
five subjects, none correct. Only inspection of annotated frames
distinguishes the two cases. Those false positives were eliminated by
raising the confidence threshold and narrowing the camera field of view.

\begin{figure}[htbp]
\centering
\begin{subfigure}[b]{0.32\textwidth}
  \includegraphics[width=\textwidth]{detection_a.png}
  \caption{Two subjects, confidence $0.59$ and $0.42$.}
  \label{fig:test-det-a}
\end{subfigure}
\hfill
\begin{subfigure}[b]{0.32\textwidth}
  \includegraphics[width=\textwidth]{detection_b.png}
  \caption{One subject, confidence $0.48$.}
  \label{fig:test-det-b}
\end{subfigure}
\hfill
\begin{subfigure}[b]{0.32\textwidth}
  \includegraphics[width=\textwidth]{detection_c.png}
  \caption{Three subjects, confidence $0.59$, $0.47$ and $0.46$.}
  \label{fig:test-det-c}
\end{subfigure}
\caption{Detector output during flight, with the edge node's log above each
view. Bounding boxes coincide with the human models rather than with ground
texture. The log records every frame, including those with no detection, so
that a stalled node is distinguishable from an empty scene, and reports
inference duration per frame.}
\label{fig:test-detection}
\end{figure}

Figure~\ref{fig:test-detection} shows three moments from a patrol flight.
Between one and four of the five subjects were detected per frame, at
confidences of \(0.42\)--\(0.59\). The shortfall against the known count of
five is attributable to scene geometry rather than to detector deficiency:
the models stand approximately one metre apart, so from the oblique camera
view the nearer figures partially occlude those behind them, and a partially
occluded subject at approximately \(53\) pixels in height falls below the
confidence threshold. Because the number of subjects is known exactly, any
count exceeding five would be a false positive without the need to
adjudicate individual boxes; none was observed in this configuration.

Inference duration was measured around the model invocation alone, excluding
message reception, format conversion and publication, so that the figure
remains comparable between back-ends and against standalone execution. Over
\(480\) consecutive frames of one patrol flight the mean was
\(1{,}093\)\,ms, with a minimum of \(1{,}025\)\,ms. The distribution is tight
apart from a single outlier at \(6{,}949\)\,ms, one frame in \(480\),
attributable to contention for the processor from another task rather than to
the model itself.

The Pi therefore completes rather less than one detection per second against a
camera delivering five frames per second, and discards approximately \(80\%\)
of what it receives --- which is why the camera subscription uses
\texttt{BEST\_EFFORT} with queue depth one, keeping the newest frame rather
than accumulating a backlog of stale ones.

% ---------------------------------------------------------------------
\subsection{Result 3: validating the measurement chain}
\label{sec:test-chain}

The two results above depend on the camera actually reaching the edge node
at the rate the configuration specifies. That assumption failed, and the
manner of its failure is reported here because it determines what the other
measurements are worth.

The symptom was that the edge node received imagery at \(0.19\)\,Hz where
the camera was configured for \(5\)\,Hz, a shortfall of approximately
\(26\times\). The visible consequences appeared three layers above the
cause: the detector reported no humans in \(93\) of \(94\) frames, and the
patrol mission timed out at every waypoint. Both invite explanation in terms
of detector settings or flight control, and neither has anything to do with
either --- the aircraft flew correctly and the detector functioned
correctly, but the scene was being observed through roughly one frame every
five seconds.

\begin{table}[htbp]
\centering
\caption{Eliminating candidate causes of the frame-rate shortfall. Each row
is a measurement that excludes one layer.}
\label{tab:test-elimination}
\begin{tabular}{@{}p{38mm}p{38mm}p{40mm}@{}}
\toprule
Hypothesis & Measurement & Outcome \\
\midrule
Simulation running slowly
  & Real-time factor from \texttt{/clock}
  & \(0.998\); simulated time tracks wall clock \\[3pt]
Gazebo rendering slowly
  & Host-side subscriber on the camera topic
  & \(5.00\)\,Hz, gaps \(0.19\)--\(0.21\)\,s; source correct \\[3pt]
Link negotiated at $100$\,Mbps
  & \texttt{ethtool} on the edge node
  & $1000$\,Mb/s full duplex \\[3pt]
Fragments lost in flight
  & Interface counters, both machines
  & Zero errors, zero drops \\[3pt]
Receive buffers too small
  & \texttt{net.core.rmem\_max}
  & $512$\,MB, already raised \\[3pt]
\textbf{Send buffers too small}
  & \texttt{net.core.wmem\_max}
  & \(\mathbf{208}\)\,\textbf{kB}, left at the default \\
\bottomrule
\end{tabular}
\end{table}

The DDS profile requests a \(16\)\,MB send buffer. The kernel caps any such
request at \texttt{net.core.wmem\_max}, and performs the reduction silently.
A single frame is therefore approximately \(13\times\) larger than the
buffer available to transmit it: the publisher filled that buffer
immediately and discarded the remaining fragments at the socket layer,
before they reached the network interface. Comparing production with
transmission confirmed this --- the host generated \(414\)\,MB of imagery
over thirty seconds and placed \(10\)\,MB of it on the wire, \(2\)\,Mbps
against the \(110\)\,Mbps the configuration implies. Raising the parameter
on both machines and restarting the stack, which is required because
Fast~DDS applies socket options only at participant creation, restored
delivery to \(5.000\)\,Hz with an inter-arrival standard deviation of
\(4.5\)\,ms: an exact match to the source.

Two properties of this fault are worth recording. First, no counter
registers it. Loss at the socket layer is invisible to interface
statistics, and \texttt{tx\_dropped} on the sending interface and
\texttt{rx\_errors} on the receiving one both remained at zero throughout,
because the discarded fragments were never transmitted. A diagnosis founded
on interface counters would have concluded that the network was healthy,
which it was. Second, the two directions are governed by separate kernel
parameters. The receive path had been corrected during earlier work and the
transmit path had not, because the failure that prompted the original
correction pointed only at the receiving side. The pre-flight checks in
\texttt{scripts/netns/run\_hitl.sh} now verify both parameters on both
machines before the stack starts.

% ---------------------------------------------------------------------
\subsection{Limitations}
\label{sec:test-limitations}

Clock synchronisation between the two machines has not been established.
End-to-end latency would be computed from timestamps taken on different
hosts, and without synchronisation such values are not meaningful while
remaining plausible in magnitude. End-to-end latency is therefore not
reported, rather than reported with a caveat.

The edge configuration is characterised; the corresponding ground
configuration, in which compressed imagery crosses the simulated channel and
inference is performed at the ground station, remains to be measured. All
measurements were taken with the aircraft close to the ground station,
giving zero packet loss on both paths, so the reliability policy applied to
detection messages is untested under the loss conditions it exists to
handle. The testbed supports a single physical edge node; multi-node
operation has not been evaluated.

% =====================================================================
%  IMAGE CHECKLIST — save into docs/report/figures/
%
%  detection_a.png / detection_b.png / detection_c.png
%      The three screenshots already captured: the detector log above and
%      the rqt_image_view annotated frame below, showing 2, 1 and 3
%      detections respectively.
%
%      CROP EACH ONE before including. At three-to-a-row the full desktop
%      screenshot will be illegible. Keep only:
%        - the last ~8 log lines (enough to show the count and the
%          inference time), and
%        - the image pane itself, without the rqt toolbar and topic
%          selector.
%      Crop all three to the same aspect ratio so the row lines up.
%
%      If the log text is still too small at 0.32\textwidth, drop the log
%      from the figure entirely and show only the three annotated views —
%      the inference figures are already stated in the prose.
%
%  OPTIONAL, if a fourth image is wanted:
%      testbed_hardware.jpg — a photograph of the laptop, the USB gigabit
%      adapter, the single Ethernet cable and the Pi 4B. This is the only
%      image that shows the work is hardware-in-the-loop rather than pure
%      simulation. It would go in Section \ref{sec:test-env}.
% =====================================================================
